\documentclass[12pt,openany]{book}
\usepackage{notes-preamble}

\allowdisplaybreaks
\emergencystretch=3em
\setcounter{secnumdepth}{2}
\setcounter{tocdepth}{2}

% Project-specific number-theory notation.
\newcommand{\Z}{\mathbb{Z}}
\newcommand{\Nzero}{\mathbb{N}_0}
\DeclareMathOperator{\ord}{ord}
\DeclareMathOperator{\lcm}{lcm}
\newcommand{\mobius}{\mu}
\newcommand{\totient}{\varphi}
\newcommand{\legendre}[2]{\left(\frac{#1}{#2}\right)}
\newcommand{\jacobi}[2]{\left(\frac{#1}{#2}\right)}
\newcommand{\overbar}[1]{\overline{#1}}

% Editorial source markers occupy no extra vertical space, preventing a marker
% from creating an otherwise blank page at a chapter boundary.
\newcommand{\sourcepage}[1]{\par\nobreak\nointerlineskip\hbox to \linewidth{\hfill\smash{\footnotesize\textcolor{gray}{Source PDF p.~#1}}}}
\newcommand{\sourcepages}[1]{\par\nobreak\nointerlineskip\hbox to \linewidth{\hfill\smash{\footnotesize\textcolor{gray}{Source PDF pp.~#1}}}}

\title{MAT4003 Number Theory}
\author{}
\date{15 July 2026}


\begin{document}
	\pagenumbering{gobble}
	\maketitle
	\allowdisplaybreaks
	\pagenumbering{arabic}
	\begin{center}
		\Large\textbf{Preface}
    \end{center}~\
    This is the note of the course \textit{Number Theory} taught in 2021 Fall by Prof. Mario Huang at The Chinese University of Hong Kong, Shenzhen. Most of the content is based on the presentation of the instructor and written down in my understanding, so there are unavoidable mistakes and typos in the note. The note is further transcribed and edited with the help from ChatGPT. Topics of the note include: basis of number theory, arithmetic functions, primitive roots, quadratic residues, continued fractions, diophantine equations, and roots of unity.

	\tableofcontents

\setcounter{chapter}{-1}
\chapter{Preliminaries}

\begin{thm}[principle of mathematical induction]
Let $P(n)$ be a statement for $n\in\Nzero$.  If $P(0)$ is true and, for every $k\in\Nzero$, $P(k)$ implies $P(k+1)$, then $P(n)$ is true for every $n\in\Nzero$.
\end{thm}

\begin{thm}[strong induction]
Let $P(n)$ be a statement for $n\in\Nzero$.  If $P(0)$ is true and, for every $k\in\Nzero$, the truth of $P(0),\ldots,P(k)$ implies $P(k+1)$, then $P(n)$ is true for every $n\in\Nzero$.
\end{thm}

\begin{thm}[well-ordering principle]
Every nonempty subset of $\Nzero$ has a least element.
\end{thm}

\begin{thm}[pigeonhole principle]
If $S$ and $T$ are finite sets and $f:S\to T$ satisfies $\abs*{S}>\abs*{T}$, then $f$ is not injective.  Equivalently, if more than $\abs*{T}$ objects are placed into $\abs*{T}$ boxes, some box contains at least two objects.
\end{thm}

\chapter{Divisibility and Primes}

\begin{de}
Let $a,b\in\Z$ with $a\ne0$.  We say that $a$ divides $b$, and write $a\mid b$, if $b=ak$ for some $k\in\Z$.  Then $a$ is a divisor of $b$, and $b$ is a multiple of $a$.
\end{de}

\begin{thm}
For integers $a,b,c,x,y$:
\begin{enumerate}[label=(\alph*)]
\item $a\mid a$;
\item if $a\mid b$ and $b\mid c$, then $a\mid c$;
\item if $a\mid b$ and $a\mid c$, then $a\mid bx+cy$;
\item if $a\mid b$ and $b\ne0$, then $\abs a\le\abs b$.
\end{enumerate}
\end{thm}

\begin{thm}[division algorithm, positive divisor]
Let $a\in\N$ and $b\in\Z$.  There exist unique integers $q,r$ such that
\[
 b=aq+r,\qquad 0\le r<a.
\]
\end{thm}
\begin{proof}
Consider
\[
 S=\{b-aq:q\in\Z,\ b-aq\ge0\}.
\]
The set is nonempty, because $q$ can be chosen sufficiently negative.  By the well-ordering principle, $S$ has a least element $r=b-aq$.  If $r\ge a$, then $r-a=b-a(q+1)$ is a smaller nonnegative element of $S$, a contradiction.  Hence $0\le r<a$.

If $b=aq_1+r_1=aq_2+r_2$ with $0\le r_1,r_2<a$, then
$a(q_1-q_2)=r_2-r_1$.  The right side has absolute value less than $a$, so it is the only multiple of $a$ in that range, namely $0$.  Thus $r_1=r_2$ and $q_1=q_2$.
\end{proof}

\begin{thm}[division algorithm]
Let $a,b\in\Z$ with $a\ne0$.  There exist unique integers $q,r$ such that
\[
 b=aq+r,\qquad 0\le r<\abs a.
\]
\end{thm}
\begin{proof}
Apply Theorem 1.2 with divisor $\abs a$.  If $a>0$, the result is immediate.  If $a<0$, replace the quotient for $\abs a$ by its negative.  Uniqueness follows exactly as in Theorem 1.2.
\end{proof}

\begin{de}
A positive integer $d$ is a common divisor of $a$ and $b$ if $d\mid a$ and $d\mid b$.  For integers $a,b$, not both zero, the greatest common divisor $\gcd(a,b)$ is the largest positive common divisor of $a$ and $b$.
\end{de}

\begin{thm}[B\'ezout's lemma]
Let $a,b\in\Z$, not both zero.  There exist $x_0,y_0\in\Z$ such that
\[
 \gcd(a,b)=ax_0+by_0.
\]
\end{thm}
\begin{proof}
Let
\[
 S=\{ax+by:x,y\in\Z,\ ax+by>0\}.
\]
This set is nonempty, so let $d=ax_0+by_0$ be its least element.  Every common divisor of $a$ and $b$ divides $d$.

Divide $a$ by $d$: $a=dq+r$ with $0\le r<d$.  Then
\[
 r=a-dq=a(1-qx_0)+b(-qy_0)
\]
is an integer linear combination of $a$ and $b$.  If $r>0$, then $r\in S$ and $r<d$, a contradiction; hence $r=0$ and $d\mid a$.  Similarly $d\mid b$.  Therefore $d$ is a common divisor divisible by every common divisor, so $d=\gcd(a,b)$.
\end{proof}

\begin{cor}
If $d=\gcd(a,b)$, then
\[
 \{ax+by:x,y\in\Z\}=\{kd:k\in\Z\}.
\]
\end{cor}
\begin{proof}
B\'ezout's lemma gives $d=ax_0+by_0$, so every multiple of $d$ is an integer linear combination of $a$ and $b$.  Conversely, $d\mid a$ and $d\mid b$, so $d$ divides every such linear combination.
\end{proof}

\begin{thm}
For $d\in\N$ and integers $a,b$, not both zero, the following are equivalent:
\begin{enumerate}[label=(\alph*)]
\item $d=\gcd(a,b)$;
\item $d$ is the least positive integer of the form $ax+by$;
\item $d$ is a common divisor of $a$ and $b$, and every common divisor of $a$ and $b$ divides $d$.
\end{enumerate}
\end{thm}
\begin{proof}
The equivalence of (a) and (b) is Theorem 1.4.  If (b) holds, the proof of B\'ezout's lemma shows that $d\mid a$ and $d\mid b$, while every common divisor divides $d$, giving (c).  Condition (c) says precisely that $d$ is the greatest positive common divisor, giving (a).
\end{proof}

\begin{de}
Integers $a$ and $b$ are coprime, or relatively prime, if $\gcd(a,b)=1$.
\end{de}

\begin{cor}
Integers $a$ and $b$ are coprime if and only if there exist $x_0,y_0\in\Z$ such that
\[
 ax_0+by_0=1.
\]
\end{cor}

\begin{thm}
If $\gcd(a,b)=d$, then
\[
 \gcd\paren*{\frac ad,\frac bd}=1.
\]
\end{thm}
\begin{proof}
Divide a B\'ezout identity $ax_0+by_0=d$ by $d$.
\end{proof}

\begin{thm}[Euclid's lemma]
If $c\mid ab$ and $\gcd(a,c)=1$, then $c\mid b$.
\end{thm}
\begin{proof}
Choose $x_0,y_0\in\Z$ with $ax_0+cy_0=1$.  Multiplying by $b$ gives
$b=abx_0+bcy_0$, and both terms on the right are divisible by $c$.
\end{proof}

\begin{thm}
If $\gcd(a,c)=\gcd(b,c)=1$, then $\gcd(ab,c)=1$.
\end{thm}
\begin{proof}
Choose $ax+cy=1$ and $bu+cv=1$.  Multiplication gives
\[
 1=ab(xu)+c(ayv+bux+cyv),
\]
so $1$ is an integer linear combination of $ab$ and $c$.
\end{proof}

\begin{lem}
If $b=aq+r$, then
\[
 \gcd(a,b)=\gcd(a,r).
\]
\end{lem}
\begin{proof}
An integer divides both $a$ and $b$ if and only if it divides both $a$ and $b-aq=r$.  Thus the two pairs have exactly the same common divisors.
\end{proof}

\begin{thm}[Euclidean algorithm]
Let $a,b\in\N$ with $b>a$.  Repeated division gives
\begin{align*}
 b&=aq_1+r_1, &0\le r_1<a,\\
 a&=r_1q_2+r_2, &0\le r_2<r_1,\\
 &\ \vdots\\
 r_{n-2}&=r_{n-1}q_n+r_n, &0\le r_n<r_{n-1},\\
 r_{n-1}&=r_nq_{n+1}.
\end{align*}
Then $\gcd(a,b)=r_n$, the last nonzero remainder.
\end{thm}
\begin{proof}
The positive remainders form a strictly decreasing sequence, so the process terminates.  Repeated use of Lemma 1.11 gives
\[
 \gcd(a,b)=\gcd(a,r_1)=\cdots=\gcd(r_{n-1},r_n)=r_n.
\]
\end{proof}

\begin{app}[extended Euclidean algorithm]
Track coefficients $r_i=ax_i+by_i$ alongside the Euclidean algorithm by starting with
\[
 r_{-1}=b,\ (x_{-1},y_{-1})=(0,1),\qquad
 r_0=a,\ (x_0,y_0)=(1,0),
\]
and using
\[
 r_i=r_{i-2}-q_i r_{i-1},\quad
 x_i=x_{i-2}-q_i x_{i-1},\quad
 y_i=y_{i-2}-q_i y_{i-1}.
\]
For example,
\[
 \gcd(123,45)=3=45\cdot 11-123\cdot 4.
\]
\end{app}

\begin{de}
A Diophantine equation is an equation whose solutions are required to be integers.
\end{de}

\begin{thm}[linear Diophantine equation]
Let $d=\gcd(a,b)$.  The equation
\[
 ax+by=c
\]
has an integer solution if and only if $d\mid c$.  If $(x_0,y_0)$ is one solution, then all solutions are
\[
 x=x_0+k\frac bd,\qquad y=y_0-k\frac ad,\qquad k\in\Z.
\]
\end{thm}
\begin{proof}
By Corollary 1.5, $c$ is an integer linear combination of $a$ and $b$ exactly when $d\mid c$.

If $(x,y)$ and $(x_0,y_0)$ are solutions, then
$a(x-x_0)=-b(y-y_0)$.  Dividing by $d$ and using
$\gcd(a/d,b/d)=1$, Euclid's lemma gives $b/d\mid x-x_0$.  Thus $x-x_0=k(b/d)$, and substitution gives $y-y_0=-k(a/d)$.  The displayed family is easily checked directly.
\end{proof}

\begin{de}
An integer $p>1$ is prime if its only positive divisors are $1$ and $p$.  An integer greater than $1$ that is not prime is composite.
\end{de}

\begin{thm}
If $p$ is prime and $p\mid ab$, then $p\mid a$ or $p\mid b$.
\end{thm}
\begin{proof}
If $p\nmid a$, then $\gcd(a,p)=1$, so Theorem 1.9 gives $p\mid b$.
\end{proof}

\begin{thm}
If $p$ is prime and $p\mid a_1a_2\cdots a_n$, then $p\mid a_i$ for some $i$.
\end{thm}
\begin{proof}
Induct on $n$, using Theorem 1.14 at the induction step.
\end{proof}

\begin{thm}[fundamental theorem of arithmetic]
Every integer $n\ge2$ is a product of primes.  This factorization is unique up to the order of the prime factors.
\end{thm}
\begin{proof}
For existence, use strong induction.  If $n$ is prime, there is nothing to prove.  If $n$ is composite, write $n=rs$ with $1<r,s<n$ and apply the induction hypothesis to $r$ and $s$.

For uniqueness, suppose that a smallest counterexample has two factorizations
\[
 n=p_1\cdots p_r=q_1\cdots q_s.
\]
Theorem 1.15 implies that $p_1=q_j$ for some $j$; reorder so that $j=1$ and cancel this common factor.  Then $n/p_1$ is a smaller counterexample, a contradiction.
\end{proof}

\begin{de}
For nonzero integers $a,b$, the least common multiple $\lcm(a,b)$ is the least positive integer divisible by both $a$ and $b$.  The definition extends analogously to finitely many nonzero integers.
\end{de}

\begin{thm}[monic rational-root criterion]
Let
\[
 f(x)=x^n+c_{n-1}x^{n-1}+\cdots+c_1x+c_0\in\Z[x].
\]
Every rational root of $f$ is an integer.  Consequently, every real root is either an integer or irrational.
\end{thm}
\begin{proof}
Write a rational root in lowest terms as $p/q$, with $q>0$.  Multiplying $f(p/q)=0$ by $q^n$ gives
\[
 p^n+c_{n-1}p^{n-1}q+\cdots+c_0q^n=0.
\]
Thus $q\mid p^n$.  Since $\gcd(p,q)=1$, we have $q=1$.
\end{proof}

\begin{thm}
There are infinitely many primes.
\end{thm}
\begin{proof}
If $p_1,\ldots,p_n$ were all primes, the integer $N=p_1\cdots p_n+1$ would have a prime divisor.  None of the $p_i$ divides $N$, a contradiction.
\end{proof}

\begin{thm}
For every $k\in\N$, there are $k$ consecutive composite integers.
\end{thm}
\begin{proof}
The integers
\[
 (k+1)!+2,(k+1)!+3,\ldots,(k+1)!+(k+1)
\]
are consecutive, and the $i$th one in the list is divisible by $i$ for $2\le i\le k+1$.
\end{proof}

\chapter{Theory of Congruences}

\begin{de}
Let $a,b,m\in\Z$ with $m\ne0$.  We say that $a$ is congruent to $b$ modulo $m$, and write
\[
 a\equiv b\pmod m,
\]
if $m\mid(a-b)$.  The integer $m$ is the modulus.
\end{de}

\begin{thm}
Congruence modulo $m$ is an equivalence relation: for all integers $a,b,c$,
\[
 a\equiv a\pmod m,
\]
$a\equiv b\pmod m$ implies $b\equiv a\pmod m$, and
$a\equiv b\pmod m$, $b\equiv c\pmod m$ imply $a\equiv c\pmod m$.
\end{thm}

\begin{thm}
If $a\equiv b\pmod m$ and $c\equiv d\pmod m$, then
\[
 a+c\equiv b+d\pmod m,
 \qquad
 ac\equiv bd\pmod m.
\]
\end{thm}

\begin{thm}[cancellation]
Let $d=\gcd(c,m)$, with $m\ne0$.  Then
\[
 ca\equiv cb\pmod m
 \quad\Longleftrightarrow\quad
 a\equiv b\pmod{m/d}.
\]
\end{thm}
\begin{proof}
The left side is $m\mid c(a-b)$.  Dividing by $d$ gives
\[
 \frac md\mid \frac cd(a-b).
\]
Because $\gcd(c/d,m/d)=1$, Euclid's lemma permits cancellation of $c/d$, giving $(m/d)\mid(a-b)$.  The reverse implication is immediate.
\end{proof}

\begin{cor}
If $\gcd(c,m)=1$, then
\[
 ca\equiv cb\pmod m\quad\Longleftrightarrow\quad a\equiv b\pmod m.
\]
\end{cor}

\begin{de}
A set $\{x_1,\ldots,x_m\}$ of integers is a complete residue system modulo $m\ge1$ if no two elements are congruent modulo $m$.
\end{de}

\begin{de}
A congruence
\[
 ax\equiv b\pmod m
\]
with unknown $x$ is a linear congruence in one variable.  A solution means a congruence class modulo $m$.
\end{de}

\begin{thm}[linear congruences]
Let $m\ge1$ and $d=\gcd(a,m)$.  The congruence
\[
 ax\equiv b\pmod m
\]
has a solution if and only if $d\mid b$.  When it is solvable, it has exactly $d$ incongruent solutions modulo $m$.
\end{thm}
\begin{proof}
The congruence is equivalent to the Diophantine equation
\[
 ax-my=b.
\]
By Theorem 1.13 this equation is solvable exactly when $d\mid b$.

Fix one solution $x_0$.  The integer solutions have
\[
 x=x_0+t\frac md,\qquad t\in\Z.
\]
Modulo $m$, the values $t=0,1,\ldots,d-1$ are pairwise incongruent, and every other value of $t$ is congruent to one of them.  Hence there are exactly $d$ solution classes.
\end{proof}

\begin{de}
If $\gcd(a,m)=1$, Theorem 2.4 gives a unique $x\pmod m$ satisfying $ax\equiv1\pmod m$.  This class is the inverse of $a$ modulo $m$ and is denoted $a^{-1}$.
\end{de}

\begin{thm}[Chinese remainder theorem]
Let $m_1,\ldots,m_r$ be pairwise coprime positive integers and let $a_1,\ldots,a_r\in\Z$.  The system
\[
 x\equiv a_i\pmod{m_i},\qquad 1\le i\le r,
\]
has a solution, unique modulo $M=m_1\cdots m_r$.
\end{thm}

\begin{lem}
If $m_1,\ldots,m_r$ are pairwise coprime and $x\equiv y\pmod{m_i}$ for every $i$, then
\[
 x\equiv y\pmod{m_1\cdots m_r}.
\]
\end{lem}
\begin{proof}
Induct on $r$.  If $m_1\cdots m_{r-1}\mid x-y$ and $m_r\mid x-y$, pairwise coprimality gives
$\gcd(m_1\cdots m_{r-1},m_r)=1$, so their product divides $x-y$.
\end{proof}

\begin{proof}[Proof of the Chinese remainder theorem]
Set $M_i=M/m_i$.  Since $\gcd(M_i,m_i)=1$, choose $u_i$ with
\[
 M_i u_i\equiv1\pmod{m_i}.
\]
Then
\[
 x_0=\sum_{i=1}^r a_iM_i u_i
\]
satisfies every congruence: modulo $m_j$, all terms vanish except the $j$th, which is congruent to $a_j$.  If $x$ and $y$ are two solutions, Lemma 2.6 gives $x\equiv y\pmod M$.
\end{proof}

\begin{thm}[Fermat's little theorem]
If $p$ is prime and $p\nmid a$, then
\[
 a^{p-1}\equiv1\pmod p.
\]
Equivalently, for every integer $a$,
\[
 a^p\equiv a\pmod p.
\]
\end{thm}

\begin{de}
Euler's totient function is
\[
 \totient(m)=\abs*{\{1\le k\le m:\gcd(k,m)=1\}}.
\]
\end{de}

\begin{de}
A set $\{x_1,\ldots,x_{\totient(m)}\}$ is a reduced residue system modulo $m$ if each $x_i$ is coprime to $m$ and no two are congruent modulo $m$.
\end{de}

\begin{remark}
If $a\equiv a'\pmod m$, then $\gcd(a,m)=\gcd(a',m)$.  Thus a reduced residue system is obtained by choosing one representative from every invertible congruence class modulo $m$.
\end{remark}

\begin{lem}
If $\gcd(a,m)=1$ and $x_1,\ldots,x_{\totient(m)}$ is a reduced residue system modulo $m$, then
\[
 ax_1,\ldots,ax_{\totient(m)}
\]
is also a reduced residue system modulo $m$.
\end{lem}
\begin{proof}
Each product is coprime to $m$.  If $ax_i\equiv ax_j\pmod m$, cancellation gives $x_i\equiv x_j\pmod m$, hence $i=j$.
\end{proof}

\begin{thm}[Euler's theorem]
If $\gcd(a,m)=1$, then
\[
 a^{\totient(m)}\equiv1\pmod m.
\]
\end{thm}
\begin{proof}
Let $x_1,\ldots,x_{\totient(m)}$ be a reduced residue system.  By Lemma 2.9, the two products
\[
 x_1\cdots x_{\totient(m)}
 \quad\text{and}\quad
 (ax_1)\cdots(ax_{\totient(m)})
\]
are congruent modulo $m$.  Their common factor $x_1\cdots x_{\totient(m)}$ is invertible modulo $m$, so cancellation gives the result.
\end{proof}

\begin{proof}[Proof of Theorem 2.7]
For prime $p$, $\totient(p)=p-1$, so the first assertion is Euler's theorem.  If $p\mid a$, then $a^p\equiv a\equiv0\pmod p$; otherwise multiply $a^{p-1}\equiv1$ by $a$.
\end{proof}

\begin{thm}[Freshman's dream modulo a prime]
For integers $x,y$ and prime $p$,
\[
 (x+y)^p\equiv x^p+y^p\pmod p.
\]
\end{thm}
\begin{proof}
By Fermat's little theorem, $(x+y)^p\equiv x+y$, $x^p\equiv x$, and $y^p\equiv y$ modulo $p$.
\end{proof}

\begin{lem}
If $p$ is prime, then
\[
 x^2\equiv1\pmod p
 \quad\Longleftrightarrow\quad
 x\equiv\pm1\pmod p.
\]
\end{lem}
\begin{proof}
The congruence is equivalent to $p\mid(x-1)(x+1)$, so primality gives the two possibilities.
\end{proof}

\begin{thm}[Wilson's theorem]
If $p$ is prime, then
\[
 (p-1)!\equiv-1\pmod p.
\]
\end{thm}
\begin{proof}
The case $p=2$ is immediate.  For odd $p$, pair every element of the reduced residue system $1,2,\ldots,p-1$ with its inverse.  By Lemma 2.12, the only self-inverse classes are $1$ and $-1$.  Every other pair contributes $1$, so the product is $-1$ modulo $p$.
\end{proof}

\begin{thm}[converse to Wilson]
Let $n\ge2$.  If
\[
 (n-1)!\equiv-1\pmod n,
\]
then $n$ is prime.
\end{thm}
\begin{proof}
Suppose $n$ is composite.  The case $n=4$ gives $(n-1)!=6\not\equiv-1\pmod4$.  For $n\ge6$, write $n=ab$ with $2\le a\le b\le n-1$.  If $a<b$, both $a$ and $b$ occur among the factors of $(n-1)!$, so $n\mid(n-1)!$.  If $a=b$, then $n=a^2$ with $a\ge3$; the factors $a$ and $2a$ both occur, so again $a^2\mid(n-1)!$.  Either conclusion contradicts $(n-1)!\equiv-1\pmod n$.
\end{proof}

\begin{thm}
Let $p$ be prime.  The congruence
\[
 x^2\equiv-1\pmod p
\]
has a solution if and only if $p=2$ or $p\equiv1\pmod4$.
\end{thm}
\begin{proof}
The case $p=2$ is immediate.  Let $p$ be odd.  If $x^2\equiv-1\pmod p$, then
\[
 1\equiv x^{p-1}=(x^2)^{(p-1)/2}\equiv(-1)^{(p-1)/2}\pmod p,
\]
so $(p-1)/2$ is even and $p\equiv1\pmod4$.

Conversely, suppose $p\equiv1\pmod4$ and no solution exists.  Pair each nonzero residue $u$ with $-u^{-1}$.  The absence of a solution to $u^2\equiv-1$ ensures that the two members of a pair are distinct, and each pair has product $-1$.  Hence
\[
 (p-1)!\equiv(-1)^{(p-1)/2}\equiv1\pmod p,
\]
contradicting Wilson's theorem.
\end{proof}

\section{Primality tests}

\begin{de}
A primality test is an algorithm that determines whether a given positive integer is prime.
\end{de}

\begin{remark}[trial division]
To test $n>1$ deterministically, it is enough to check divisibility by primes not exceeding $\sqrt n$.  The test is exact but slow for large inputs.
\end{remark}

\begin{app}[Fermat primality test]
For an odd input $n$, choose random bases $a\in\{2,\ldots,n-2\}$.  If
$a^{n-1}\not\equiv1\pmod n$, then $n$ is composite.  If many independent bases pass, return ``probably prime.''
\end{app}

\begin{de}
Let $n$ be composite and $1\le a<n$.  If $a^{n-1}\equiv1\pmod n$, then $n$ is a Fermat pseudoprime to base $a$, and $a$ is a Fermat liar.  Otherwise $a$ is a Fermat witness.
\end{de}

\begin{de}
A composite integer $n$ is a Carmichael number if
\[
 a^{n-1}\equiv1\pmod n
\]
for every $a$ coprime to $n$.
\end{de}
\begin{remark}
Carmichael numbers show that the Fermat test can fail for every coprime base.  Infinitely many Carmichael numbers exist.
\end{remark}

\begin{app}[Miller--Rabin primality test]
Let $n>2$ be odd and write
\[
 n-1=2^s d,\qquad d\text{ odd}.
\]
For a chosen base $a\in\{2,\ldots,n-2\}$:
\begin{enumerate}
\item if $\gcd(a,n)>1$, return ``composite'';
\item compute $x=a^d\bmod n$; if $x\equiv1$ or $x\equiv-1$, the base passes;
\item otherwise square repeatedly, at most $s-1$ times.  If a square is $-1$, the base passes; if a square is $1$ before $-1$, or no $-1$ occurs, return ``composite.''
\end{enumerate}
A base that passes is a strong liar; a base that exposes compositeness is a strong witness.
\end{app}

\begin{thm}[Rabin's bound]
Let $n>4$ be odd and composite.  At most one quarter of the residues $a\in\{1,\ldots,n-1\}$ are strong liars for the Miller--Rabin test.  Thus independent random trials return ``probably prime'' for a composite input with probability at most $4^{-k}$ after $k$ rounds.
\end{thm}
\begin{proof}
Write
\[
 n=\prod_{j=1}^r p_j^{e_j},\qquad n-1=2^s d,
\]
where $d$ is odd and the $p_j$ are distinct odd primes.  A strong liar must be a unit.  The group
$G_j=(\Z/p_j^{e_j}\Z)^\times$ is cyclic of order
\[
 \totient(p_j^{e_j})=2^{s_j}u_j,
 \qquad u_j\text{ odd}.
\]
Put $h_j=\gcd(d,u_j)$ and
$\ell=\min(s,s_1,\ldots,s_r)$.  In a cyclic group of order $2^{s_j}u_j$, the equation $x^d=1$ has $h_j$ solutions, while
$x^{2^i d}=-1$ has $2^i h_j$ solutions exactly when $i<s_j$.  The Chinese remainder theorem therefore gives the exact number
\[
 L=\paren*{\prod_{j=1}^r h_j}
 \left(1+\sum_{i=0}^{\ell-1}2^{ir}\right)
\]
of strong liars.

It remains to bound this expression.  If $r\ge3$, then
\[
 1+\sum_{i=0}^{\ell-1}2^{ir}\le 2^{r\ell-2},
\]
so $L\le\totient(n)/4\le(n-1)/4$.

Suppose $r=2$.  Here
\[
 1+\sum_{i=0}^{\ell-1}2^{2i}=\frac{4^\ell+2}{3}.
\]
If $s_1+s_2\ge2\ell+1$, then
\[
 \frac{L}{\totient(n)}
 \le \frac{4^\ell+2}{3\cdot 2^{2\ell+1}}
 \le\frac14.
\]
The remaining case is $s_1=s_2=\ell$.  Both $h_j$ cannot equal $u_j$: if they did, the exponents $e_j$ would have to be $1$, because $p_j\nmid d$, and the divisibilities $u_1\mid d$ and $u_2\mid d$ would imply $u_1\mid u_2$ and $u_2\mid u_1$, hence $p_1=p_2$, impossible.  Thus one factor $h_j/u_j$ is at most $1/3$, and since $(4^\ell+2)/(3\cdot4^\ell)\le1/2$, we again obtain $L\le\totient(n)/6<(n-1)/4$.

Finally, if $r=1$, then $n=p^e$ with $e\ge2$.  Since $p\nmid n-1$, the factor $p^{e-1}$ in $u_1$ does not divide $h_1$, so
\[
 L=h_1 2^\ell\le p-1\le\frac{p^e-1}{4}.
\]
This proves the bound in all cases.
\end{proof}

\section{Hensel lifting}

\begin{thm}[Hensel's lemma, simple root]
Let $p$ be prime, $k\ge1$, $f\in\Z[x]$, and $a\in\Z$.  Suppose
\[
 f(a)\equiv0\pmod{p^k},
 \qquad
 f'(a)\not\equiv0\pmod p.
\]
Then there is a unique $t\in\{0,1,\ldots,p-1\}$ such that
\[
 f(a+tp^k)\equiv0\pmod{p^{k+1}}.
\]
\end{thm}
\begin{proof}
The binomial theorem gives the first-order congruence
\[
 f(a+tp^k)\equiv f(a)+tp^k f'(a)\pmod{p^{k+1}}.
\]
After division by $p^k$, the required condition is
\[
 t f'(a)\equiv-\frac{f(a)}{p^k}\pmod p.
\]
Since $f'(a)$ is invertible modulo $p$, this linear congruence has one solution modulo $p$.
\end{proof}

\begin{remark}
Every simple root modulo $p^k$ lifts uniquely to a root modulo $p^{k+1}$.
\end{remark}

\begin{thm}[singular-root alternatives]
Let $p$ be prime, $k\ge1$, $f\in\Z[x]$, and suppose
\[
 f(a)\equiv0\pmod{p^k},\qquad f'(a)\equiv0\pmod p.
\]
Then:
\begin{enumerate}[label=(\roman*)]
\item if $f(a)\equiv0\pmod{p^{k+1}}$, every $a+tp^k$, $t\in\Z$, is a root modulo $p^{k+1}$;
\item if $f(a)\not\equiv0\pmod{p^{k+1}}$, none of them is a root modulo $p^{k+1}$.
\end{enumerate}
\end{thm}
\begin{proof}
Again
\[
 f(a+tp^k)\equiv f(a)+tp^k f'(a)\equiv f(a)\pmod{p^{k+1}},
\]
because $p\mid f'(a)$.
\end{proof}

\chapter{Arithmetic Functions}

\begin{notedefinition}[3.1]
An arithmetic function is a function $f:\N\to\C$.
\end{notedefinition}
Typical examples are
\[
 \tau(n)=\sum_{d\mid n}1,
 \qquad
 \sigma(n)=\sum_{d\mid n}d,
 \qquad
 \omega(n)=\#\{p:p\mid n,\ p\text{ prime}\},
\]
and $\pi(n)$, the number of primes not exceeding $n$.

\begin{notedefinition}[3.2]
A nonzero arithmetic function $f$ is multiplicative if
\[
 f(mn)=f(m)f(n)\qquad\text{whenever }\gcd(m,n)=1.
\]
It is completely multiplicative if this equality holds for all $m,n\in\N$.
\end{notedefinition}

\begin{notetheorem}[3.1]
If
\[
 n=p_1^{\alpha_1}\cdots p_r^{\alpha_r}
\]
is the prime factorization of $n$ and $f$ is multiplicative, then
\[
 f(n)=\prod_{i=1}^r f(p_i^{\alpha_i}).
\]
\end{notetheorem}
\begin{proof}
The prime powers are pairwise coprime, so repeated use of multiplicativity gives the formula.
\end{proof}

\begin{noteremark}
Every nonzero multiplicative function satisfies $f(1)=1$.  For fixed $k\in\C$, the function $n\mapsto n^k$ is completely multiplicative.  The functions $u(n)=1$ and $r(n)=n$ are the cases $k=0$ and $k=1$.  The prime-counting function $\pi$ is not multiplicative.
\end{noteremark}

For fixed $k\in\C$, define
\[
 \sigma_k(n)=\sum_{d\mid n}d^k.
\]

\begin{noteproposition}[3.1A]
The function $\sigma_k$ is multiplicative.
\end{noteproposition}
\begin{proof}
If $\gcd(m,n)=1$, every divisor $d$ of $mn$ has a unique representation $d=d_1d_2$ with $d_1\mid m$ and $d_2\mid n$.  Therefore
\[
 \sigma_k(mn)
 =\sum_{d_1\mid m}\sum_{d_2\mid n}(d_1d_2)^k
 =\sigma_k(m)\sigma_k(n).
\]
\end{proof}

\begin{notetheorem}[3.2]
If $n=p_1^{\alpha_1}\cdots p_r^{\alpha_r}$, then
\[
 \tau(n)=\prod_{i=1}^r(\alpha_i+1),
 \qquad
 \sigma(n)=\prod_{i=1}^r\frac{p_i^{\alpha_i+1}-1}{p_i-1}.
\]
\end{notetheorem}
\begin{proof}
For a prime power,
\[
 \tau(p^\alpha)=\alpha+1,
 \qquad
 \sigma(p^\alpha)=1+p+\cdots+p^\alpha
 =\frac{p^{\alpha+1}-1}{p-1}.
\]
Now apply Theorem 3.1.
\end{proof}

\begin{notetheorem}[3.3]
Euler's totient function is multiplicative.
\end{notetheorem}
\begin{proof}
For coprime $m,n$, the Chinese remainder theorem gives a bijection
\[
 (\Z/mn\Z)^\times\longrightarrow
 (\Z/m\Z)^\times\times(\Z/n\Z)^\times,
 \qquad [x]_{mn}\longmapsto([x]_m,[x]_n).
\]
Taking cardinalities gives $\totient(mn)=\totient(m)\totient(n)$.
\end{proof}

\begin{notecorollary}[3.4]
If $n=p_1^{\alpha_1}\cdots p_r^{\alpha_r}$, then
\[
 \totient(n)=\prod_{i=1}^r(p_i^{\alpha_i}-p_i^{\alpha_i-1})
 =n\prod_{i=1}^r\left(1-\frac1{p_i}\right).
\]
\end{notecorollary}

\begin{notetheorem}[3.5 (Gauss's totient sum)]
For every $n\in\N$,
\[
 \sum_{d\mid n}\totient(d)=n.
\]
\end{notetheorem}
\begin{proof}
Partition $\{1,2,\ldots,n\}$ according to the value of $\gcd(x,n)$.  For each divisor $d\mid n$, let
\[
 S_d=\{1\le x\le n:\gcd(x,n)=d\}.
\]
Writing $x=dy$ gives a bijection from $S_d$ to the reduced residue classes modulo $n/d$, so
$\abs*{S_d}=\totient(n/d)$.  Hence
\[
 n=\sum_{d\mid n}\abs*{S_d}
 =\sum_{d\mid n}\totient(n/d)
 =\sum_{d\mid n}\totient(d).
\]
\end{proof}

\begin{notedefinition}[3.3]
A positive integer $n$ is perfect if $\sigma(n)=2n$.
\end{notedefinition}

\begin{notetheorem}[3.6 (Euclid--Euler theorem)]
An even positive integer $n$ is perfect if and only if
\[
 n=2^{m-1}(2^m-1)
\]
for some $m\in\N$ such that $2^m-1$ is prime.
\end{notetheorem}
\begin{proof}
If $2^m-1$ is prime, multiplicativity gives
\[
 \sigma\bigl(2^{m-1}(2^m-1)\bigr)
 =(2^m-1)\cdot 2^m=2n.
\]

Conversely, write an even perfect number as $n=2^k n_1$, with $n_1$ odd.  Then
\[
 (2^{k+1}-1)\sigma(n_1)=2^{k+1}n_1.
\]
Since $\gcd(2^{k+1}-1,2^{k+1})=1$, write
$n_1=(2^{k+1}-1)n_2$.  The preceding equation becomes
\[
 \sigma(n_1)=2^{k+1}n_2.
\]
If $n_2>1$, then $1,n_2,n_1$ are distinct divisors of $n_1$, so
\[
 \sigma(n_1)\ge1+n_2+n_1=1+2^{k+1}n_2,
\]
a contradiction.  Hence $n_2=1$ and $n_1=2^{k+1}-1$.  Finally,
$\sigma(n_1)=2^{k+1}=n_1+1$, so $n_1$ has no positive divisors other than $1$ and itself and is prime.  Put $m=k+1$.
\end{proof}

\begin{notetheorem}[3.7]
If $m\in\N$ and $2^m-1$ is prime, then $m$ is prime.
\end{notetheorem}
\begin{proof}
If $m=ab$ with $a,b>1$, then
\[
 2^m-1=(2^a)^b-1=(2^a-1)(1+2^a+\cdots+2^{a(b-1)}),
\]
so $2^m-1$ is composite.
\end{proof}

\section{Dirichlet convolution}

\begin{notedefinition}[3.4]
For arithmetic functions $f$ and $g$, their Dirichlet convolution is
\[
 (f*g)(n)=\sum_{d\mid n}f(d)g(n/d).
\]
Let $I$ denote the identity function
\[
 I(n)=\begin{cases}1,&n=1,\\0,&n>1.\end{cases}
\]
\end{notedefinition}

\begin{notetheorem}[3.8]
For arithmetic functions $f,g,h$:
\begin{enumerate}[label=(\alph*)]
\item $(f*g)*h=f*(g*h)$;
\item $f*g=g*f$;
\item $f*I=f$.
\end{enumerate}
\end{notetheorem}
\begin{proof}
Commutativity follows by replacing $d$ with $n/d$.  For associativity,
\[
 ((f*g)*h)(n)
 =\sum_{abc=n}f(a)g(b)h(c)
 =(f*(g*h))(n).
\]
For the identity, only the divisor $d=n$ contributes to
$\sum_{d\mid n}f(d)I(n/d)$.
\end{proof}

\begin{notedefinition}[3.5]
If $f*g=I$, then $g$ is the Dirichlet inverse of $f$, denoted $f^{-1}$.
\end{notedefinition}

\begin{notetheorem}[3.9]
An arithmetic function $f$ has a Dirichlet inverse if and only if $f(1)\ne0$.
\end{notetheorem}
\begin{proof}
If $f*g=I$, then $f(1)g(1)=1$, so $f(1)\ne0$.

Conversely, if $f(1)\ne0$, define recursively
\[
 g(1)=\frac1{f(1)},
 \qquad
 g(n)=-\frac1{f(1)}
 \sum_{\substack{d\mid n\\d>1}}f(d)g(n/d)
 \quad(n>1).
\]
Then $(f*g)(1)=1$, and for $n>1$ the defining equation makes $(f*g)(n)=0$.
\end{proof}

\begin{notetheorem}[3.10]
If $f$ and $g$ are multiplicative, then $f*g$ is multiplicative.
\end{notetheorem}
\begin{proof}
Let $\gcd(m,n)=1$.  Every divisor of $mn$ is uniquely $d_1d_2$ with $d_1\mid m$ and $d_2\mid n$.  Hence
\begin{align*}
 (f*g)(mn)
 &=\sum_{d_1\mid m}\sum_{d_2\mid n}
 f(d_1d_2)g\left(\frac m{d_1}\frac n{d_2}\right)\\
 &=\left(\sum_{d_1\mid m}f(d_1)g(m/d_1)\right)
   \left(\sum_{d_2\mid n}f(d_2)g(n/d_2)\right).
\end{align*}
\end{proof}

\begin{noteremark}
The Dirichlet inverse of a multiplicative function is multiplicative.  Indeed, define a multiplicative function $h$ on prime powers by the recursion in Theorem 3.9.  Then $f*h=I$ on prime powers and, by multiplicativity, on every integer.  Uniqueness of the inverse gives $h=f^{-1}$.
\end{noteremark}

\begin{notetheorem}[3.11]
If $f$ is multiplicative and
\[
 F(n)=\sum_{d\mid n}f(d),
\]
then $F$ is multiplicative.
\end{notetheorem}
\begin{proof}
Since $F=f*u$ and both $f$ and $u$ are multiplicative, Theorem 3.10 applies.
\end{proof}
\begin{noteremark}
Taking $f(n)=n^k$ proves that $\sigma_k$ is multiplicative.  Taking $f=\totient$ and checking prime powers gives another proof of
$\sum_{d\mid n}\totient(d)=n$.
\end{noteremark}

\begin{notedefinition}[3.6]
The M\"obius function is
\[
 \mobius(n)=
 \begin{cases}
 0,&p^2\mid n\text{ for some prime }p,\\
 (-1)^{\omega(n)},&n\text{ is squarefree}.
 \end{cases}
\]
In particular, $\mobius(1)=1$.
\end{notedefinition}

\begin{notetheorem}[3.12]
The Dirichlet inverse of $u$ is $\mobius$; equivalently,
\[
 \sum_{d\mid n}\mobius(d)=I(n).
\]
\end{notetheorem}
\begin{proof}
Both sides are multiplicative, so it suffices to check prime powers.  For $n=p^\alpha$ with $\alpha\ge1$,
\[
 \sum_{d\mid p^\alpha}\mobius(d)=\mobius(1)+\mobius(p)=1-1=0.
\]
\end{proof}

\begin{notetheorem}[3.13 (M\"obius inversion)]
If
\[
 F(n)=\sum_{d\mid n}f(d),
\]
then
\[
 f(n)=\sum_{d\mid n}\mobius(d)F(n/d)
     =\sum_{d\mid n}\mobius(n/d)F(d).
\]
\end{notetheorem}
\begin{proof}
The hypothesis is $F=f*u$.  Convolving with $\mobius=u^{-1}$ gives $f=F*\mobius$.
\end{proof}

\begin{notetheorem}[3.14 (converse form)]
If $f=\mobius*F$, then $F=f*u$; in particular, the two formulas in Theorem 3.13 are equivalent.
\end{notetheorem}

\section{The zeta function and Dirichlet series}

\begin{notedefinition}[3.7]
For $s\in\C$ with $\Re(s)>1$, the Riemann zeta function is
\[
 \zeta(s)=\sum_{n=1}^{\infty}\frac1{n^s}.
\]
\end{notedefinition}

\begin{notetheorem}[3.15 (Euler product)]
For $\Re(s)>1$,
\[
 \zeta(s)=\prod_{p\ \mathrm{prime}}\frac1{1-p^{-s}}.
\]
\end{notetheorem}
\begin{proof}
For real $s>1$, let the finite product range over $p\le N$.  Expanding its geometric series gives the sum of $n^{-s}$ over positive integers all of whose prime factors are at most $N$.  This includes every $n\le N$, so
\[
 0\le \zeta(s)-\prod_{p\le N}(1-p^{-s})^{-1}
 \le\sum_{n>N}n^{-s}\longrightarrow0.
\]
For complex $s$ with $\Re(s)>1$, the same argument applies absolutely after replacing $s$ by $\Re(s)$ in the estimates.
\end{proof}

\begin{notedefinition}[3.8]
The Dirichlet series of an arithmetic function $f$ is
\[
 D(f,s)=\sum_{n=1}^{\infty}\frac{f(n)}{n^s}
\]
wherever the series converges.
\end{notedefinition}

\begin{notetheorem}[3.16]
If $D(f,s)$ and $D(g,s)$ converge absolutely, then $D(f*g,s)$ converges absolutely and
\[
 D(f,s)D(g,s)=D(f*g,s).
\]
\end{notetheorem}
\begin{proof}
Absolute convergence gives
\[
 \sum_{a,b\ge1}\left|\frac{f(a)g(b)}{(ab)^s}\right|<\infty,
\]
so the double series may be rearranged.  Grouping terms with $ab=n$ yields
\[
 D(f,s)D(g,s)
 =\sum_{n=1}^{\infty}\frac1{n^s}\sum_{ab=n}f(a)g(b)
 =D(f*g,s).
\]
\end{proof}

\begin{notetheorem}[3.17]
For $\Re(s)>1$,
\[
 \frac1{\zeta(s)}=\sum_{n=1}^{\infty}\frac{\mobius(n)}{n^s}.
\]
\end{notetheorem}
\begin{proof}
Since $u*\mobius=I$, Theorem 3.16 gives
\[
 D(u,s)D(\mobius,s)=D(I,s)=1.
\]
\end{proof}

\section{Two natural-density applications}

\begin{noteproposition}[3.18]
The set of squarefree positive integers has natural density
\[
 \frac1{\zeta(2)}=\frac6{\pi^2}.
\]
\end{noteproposition}
\begin{proof}
The indicator of squarefreeness is
\[
 \sum_{d^2\mid n}\mobius(d).
\]
If $Q(N)$ counts squarefree integers not exceeding $N$, then
\begin{align*}
 Q(N)
 &=\sum_{d\le\sqrt N}\mobius(d)\floor{N/d^2}\\
 &=N\sum_{d\le\sqrt N}\frac{\mobius(d)}{d^2}+O(\sqrt N).
\end{align*}
After division by $N$, the error tends to zero and the finite sum tends to
$\sum_{d\ge1}\mobius(d)/d^2=1/\zeta(2)$.
\end{proof}

\begin{noteproposition}[3.19]
The proportion of ordered pairs $(a,b)$ with $1\le a,b\le N$ and $\gcd(a,b)=1$ tends to
\[
 \frac1{\zeta(2)}=\frac6{\pi^2}.
\]
\end{noteproposition}
\begin{proof}
Using $I(\gcd(a,b))=\sum_{d\mid a,\ d\mid b}\mobius(d)$, the number $A(N)$ of coprime pairs is
\[
 A(N)=\sum_{d\le N}\mobius(d)\floor{N/d}^2.
\]
Since $\floor{N/d}^2=N^2/d^2+O(N/d)$,
\[
 \frac{A(N)}{N^2}
 =\sum_{d\le N}\frac{\mobius(d)}{d^2}
 +O\left(\frac1N\sum_{d\le N}\frac1d\right)
 \longrightarrow\frac1{\zeta(2)}.
\]
\end{proof}

\chapter{Primitive Roots}

\begin{de}
Let $m\ge1$ and $\gcd(a,m)=1$.  The order of $a$ modulo $m$, denoted $\ord_m(a)$, is the least positive integer $k$ such that
\[
 a^k\equiv1\pmod m.
\]
Euler's theorem guarantees that such a $k$ exists.
\end{de}

\begin{thm}
Let $k=\ord_m(a)$.  For $h\in\N$,
\[
 a^h\equiv1\pmod m\quad\Longleftrightarrow\quad k\mid h.
\]
\end{thm}
\begin{proof}
Write $h=qk+r$ with $0\le r<k$.  Then
$a^h\equiv a^r\pmod m$.  If $a^h\equiv1$, minimality of $k$ forces $r=0$.  The converse is immediate.
\end{proof}

\begin{cor}
If $\gcd(a,m)=1$, then
\[
 \ord_m(a)\mid\totient(m).
\]
\end{cor}

\begin{thm}
If $\gcd(a,m)=1$, then for integers $i,j$,
\[
 a^i\equiv a^j\pmod m
 \quad\Longleftrightarrow\quad
 i\equiv j\pmod{\ord_m(a)}.
\]
\end{thm}
\begin{proof}
Because $a$ is invertible modulo $m$, the first congruence is equivalent to
$a^{i-j}\equiv1\pmod m$.  Apply Theorem 4.1.
\end{proof}

\begin{cor}
If $k=\ord_m(a)$, then $a,a^2,\ldots,a^k$ are pairwise incongruent modulo $m$.
\end{cor}

\begin{thm}
If $k=\ord_m(a)$ and $h\in\N$, then
\[
 \ord_m(a^h)=\frac{k}{\gcd(h,k)}.
\]
\end{thm}
\begin{proof}
Put $d=\gcd(h,k)$, $h=dh_1$, and $k=dk_1$ with $\gcd(h_1,k_1)=1$.  Then
$(a^h)^{k_1}=a^{hk_1}=a^{h_1k}\equiv1$, so the order divides $k_1$.
If $(a^h)^t\equiv1$, then $k\mid ht$, hence $k_1\mid h_1t$.  Coprimality gives $k_1\mid t$, so $k_1$ is the order.
\end{proof}

\begin{de}
A primitive root modulo $m$ is an integer $g$ satisfying
\[
 \ord_m(g)=\totient(m).
\]
\end{de}

\begin{thm}
If $g$ is a primitive root modulo $m$, then
\[
 g,g^2,\ldots,g^{\totient(m)}
\]
is a reduced residue system modulo $m$.
\end{thm}
\begin{proof}
Every power is a unit, and Corollary 4.4 shows that the $\totient(m)$ powers are incongruent.  They therefore represent all unit classes.
\end{proof}

\begin{thm}
If $m$ has a primitive root, then it has exactly
\[
 \totient(\totient(m))
\]
incongruent primitive roots modulo $m$.
\end{thm}
\begin{proof}
Let $g$ be a primitive root.  Every unit is $g^k$ for a unique $1\le k\le\totient(m)$.  By Theorem 4.5,
\[
 \ord_m(g^k)=\frac{\totient(m)}{\gcd(k,\totient(m))},
\]
which equals $\totient(m)$ exactly when $\gcd(k,\totient(m))=1$.
\end{proof}

\section{Roots of polynomial congruences}

\begin{thm}[Lagrange's root bound]
Let $p$ be prime and
\[
 f(x)=a_nx^n+\cdots+a_0\in\Z[x],
 \qquad p\nmid a_n.
\]
Then $f(x)\equiv0\pmod p$ has at most $n$ incongruent solutions modulo $p$.
\end{thm}
\begin{proof}
Induct on $n$.  If $f$ has no root, there is nothing to prove.  If $x_0$ is a root, polynomial division in $\mathbb{F}_p[x]$ gives
\[
 f(x)=(x-x_0)g(x)
\]
with $\deg g=n-1$.  Every root of $f$ is either $x_0$ or a root of $g$, so the induction hypothesis gives at most $1+(n-1)=n$ roots.
\end{proof}

\begin{cor}
If a polynomial of degree at most $n$ has more than $n$ roots modulo a prime $p$, then all its coefficients are divisible by $p$.
\end{cor}

\begin{thm}
Let $p$ be prime and $d\mid p-1$.  Then
\[
 x^d\equiv1\pmod p
\]
has exactly $d$ incongruent solutions.
\end{thm}
\begin{proof}
Write $p-1=dt$.  For every nonzero residue $x$,
\[
 x^{p-1}-1=(x^d-1)(1+x^d+\cdots+x^{d(t-1)})\equiv0\pmod p.
\]
The second factor has degree $d(t-1)=p-1-d$, so Theorem 4.8 gives it at most $p-1-d$ roots.  Among the $p-1$ nonzero residues, at least $d$ are therefore roots of $x^d-1$.  Theorem 4.8 gives the opposite bound, so there are exactly $d$.
\end{proof}

\begin{thm}
If $p$ is prime and $d\mid p-1$, then exactly $\totient(d)$ incongruent residues have order $d$ modulo $p$.
\end{thm}
\begin{proof}
Let $N_d$ be the number of residues of order $d$.  Every nonzero residue has order dividing $p-1$, so
\[
 \sum_{d\mid p-1}N_d=p-1.
\]
If $N_d>0$, choose $a$ of order $d$.  The $d$ roots of $x^d\equiv1$ are exactly
$a,a^2,\ldots,a^d$ by Theorem 4.10, and $a^j$ has order $d$ exactly when $\gcd(j,d)=1$.  Hence $N_d=\totient(d)$ whenever $N_d>0$.
Since Gauss's identity gives $\sum_{d\mid p-1}\totient(d)=p-1$, no $N_d$ can be zero, and the result follows.
\end{proof}

\begin{cor}
Every prime $p$ has primitive roots, and exactly $\totient(p-1)$ incongruent residues are primitive roots modulo $p$.
\end{cor}

\begin{remark}[finding a primitive root modulo a prime]
Factor $p-1=q_1^{e_1}\cdots q_r^{e_r}$.  A unit $a$ is primitive modulo $p$ exactly when
\[
 a^{(p-1)/q_i}\not\equiv1\pmod p
 \qquad(1\le i\le r).
\]
Testing random candidates therefore requires the distinct prime divisors of $p-1$.
\end{remark}

\section{Lifting primitive roots}

\begin{thm}
If $p$ is prime and $g$ is a primitive root modulo $p$, then at least one of $g$ and $g+p$ is a primitive root modulo $p^2$.
\end{thm}
\begin{proof}
The case $p=2$ is immediate.  Let $p$ be odd.  The order of $g$ modulo $p^2$ is a multiple of $p-1$ and divides $p(p-1)$, so it is either $p-1$ or $p(p-1)$.  If $g$ is not primitive modulo $p^2$, then
$g^{p-1}\equiv1\pmod{p^2}$.  The binomial theorem gives
\[
 (g+p)^{p-1}
 \equiv g^{p-1}+(p-1)pg^{p-2}
 \equiv1-pg^{p-2}\not\equiv1\pmod{p^2}.
\]
Thus $g+p$ has order $p(p-1)$.
\end{proof}

For a prime $p$ and nonzero integer $N$, let $v_p(N)$ be the exponent of $p$ in the prime factorization of $N$.

\begin{lem}[A lifting valuation]
If $p$ is odd and $p\mid u-1$, then
\[
 v_p(u^p-1)=v_p(u-1)+1.
\]
\end{lem}
\begin{proof}
Write $u=1+p^t c$ with $p\nmid c$.  In the binomial expansion of $u^p-1$, the first term is $p^{t+1}c$, while every remaining term is divisible by $p^{t+2}$.
\end{proof}

\begin{thm}
Let $p$ be an odd prime.  If $g$ is a primitive root modulo $p^2$, then $g$ is a primitive root modulo $p^k$ for every $k\ge2$.
\end{thm}
\begin{proof}
Proceed by induction on $k$.  The case $k=2$ is the hypothesis.  Since $g$ is primitive modulo $p^2$,
\[
 v_p(g^{p-1}-1)=1.
\]
Repeated application of the lifting valuation gives
\[
 v_p\left(g^{p^{k-2}(p-1)}-1\right)=k-1.
\]
Assume the result modulo $p^{k-1}$.  Then $\ord_{p^k}(g)$ is a multiple of
$\ord_{p^{k-1}}(g)=p^{k-2}(p-1)$ and divides
$\totient(p^k)=p^{k-1}(p-1)$.  If it were the smaller value, then $p^k$ would divide the displayed difference, contradicting its valuation.  Hence the order is $p^{k-1}(p-1)$.
\end{proof}

\begin{thm}
Let $p$ be an odd prime and let $g$ be a primitive root modulo $p^k$.  If $g$ is odd, then $g$ is a primitive root modulo $2p^k$.  If $g$ is even, then $g+p^k$ is a primitive root modulo $2p^k$.
\end{thm}
\begin{proof}
An odd $g$ has order $1$ modulo $2$ and order $\totient(p^k)$ modulo $p^k$.  By the Chinese remainder theorem, its order modulo $2p^k$ is their least common multiple, namely
$\totient(p^k)=\totient(2p^k)$.  If $g$ is even, then $g+p^k$ is odd and congruent to $g$ modulo $p^k$, so the first case applies.
\end{proof}

\begin{thm}
If $k\ge3$ and $b$ is odd, then
\[
 b^{2^{k-2}}\equiv1\pmod{2^k}.
\]
Consequently, $2^k$ has no primitive roots for $k\ge3$.
\end{thm}
\begin{proof}
For $k=3$, every odd square is $1$ modulo $8$.  If
$b^{2^{k-3}}=1+c2^{k-1}$, then squaring gives
\[
 b^{2^{k-2}}\equiv1+2c2^{k-1}\equiv1\pmod{2^k}.
\]
The exponent is $\totient(2^k)/2$, so no unit can have order $\totient(2^k)$.
\end{proof}

\begin{thm}
If $m_1,m_2\ge3$ are coprime, then $m_1m_2$ has no primitive roots.
\end{thm}
\begin{proof}
For every unit $a$ modulo $m_1m_2$, Euler's theorem gives
$a^{\totient(m_i)}\equiv1\pmod{m_i}$.  Since both totients are even,
\[
 \frac{\totient(m_1m_2)}2
 =\frac{\totient(m_1)\totient(m_2)}2
\]
is a multiple of each $\totient(m_i)$.  Hence
$a^{\totient(m_1m_2)/2}\equiv1$ modulo both $m_1$ and $m_2$, and therefore modulo their product.  No unit has full order.
\end{proof}

\begin{thm}[classification of primitive-root moduli]
An integer $m\ge2$ has a primitive root if and only if
\[
 m=2,\quad 4,\quad p^k,\quad\text{or}\quad2p^k,
\]
where $p$ is an odd prime and $k\in\N$.
\end{thm}
\begin{proof}
The moduli $2$ and $4$ are immediate.  Corollary 4.12, Theorems 4.12A and 4.13 produce primitive roots modulo every odd prime power, and Theorem 4.14 produces them modulo twice an odd prime power.

Conversely, write $m=2^\alpha n$ with $n$ odd.  If $n$ contains two distinct prime-power factors, split $m$ into two coprime factors, each at least $3$, and apply Theorem 4.16.  Thus $n=1$ or $n=p^k$.

If $\alpha\ge3$, every unit $a$ satisfies
$a^{2^{\alpha-2}}\equiv1\pmod{2^\alpha}$ by Theorem 4.15 and
$a^{\totient(n)}\equiv1\pmod n$ by Euler's theorem.  Hence
\[
 a^{\totient(m)/2}=a^{2^{\alpha-2}\totient(n)}\equiv1\pmod m,
\]
so no primitive root exists.  Finally, if $\alpha=2$ and $n=p^k>1$, the modulus is $4p^k$ and Theorem 4.16 applies to the coprime factors $4$ and $p^k$.  The only remaining forms are those listed.
\end{proof}

\section{Indices and power congruences}

\begin{de}
Suppose $m$ has primitive root $g$ and $\gcd(a,m)=1$.  The unique integer $k$ with
$1\le k\le\totient(m)$ and
\[
 g^k\equiv a\pmod m
\]
is the index, or discrete logarithm, of $a$ to the base $g$, denoted $\operatorname{ind}_g(a)$.
\end{de}

\begin{thm}
For units $a,b$ modulo $m$ and $t\in\Nzero$,
\begin{align*}
 \operatorname{ind}_g(1)&\equiv0\pmod{\totient(m)},\\
 \operatorname{ind}_g(ab)&\equiv\operatorname{ind}_g(a)+\operatorname{ind}_g(b)\pmod{\totient(m)},\\
 \operatorname{ind}_g(a^t)&\equiv t\operatorname{ind}_g(a)\pmod{\totient(m)}.
\end{align*}
\end{thm}

\begin{thm}
Let $m$ have a primitive root, let $\gcd(a,m)=1$, let $k\in\N$, and put
\[
 d=\gcd(k,\totient(m)).
\]
Then
\[
 x^k\equiv a\pmod m
\]
has a solution if and only if
\[
 a^{\totient(m)/d}\equiv1\pmod m.
\]
When solvable, it has exactly $d$ incongruent solutions modulo $m$.
\end{thm}
\begin{proof}
Any solution $x$ is a unit modulo $m$, because $x^k\equiv a$ and $a$ is a unit.  Write $a\equiv g^\ell$ and $x\equiv g^b$.  The congruence is equivalent to
\[
 kb\equiv\ell\pmod{\totient(m)}.
\]
By Theorem 2.4, this has a solution exactly when $d\mid\ell$, and then it has exactly $d$ solutions for $b$ modulo $\totient(m)$.  Also,
\[
 a^{\totient(m)/d}\equiv1
 \quad\Longleftrightarrow\quad
 \totient(m)\mid \ell\frac{\totient(m)}d
 \quad\Longleftrightarrow\quad d\mid\ell.
\]
The correspondence $b\mapsto g^b$ preserves distinct classes, completing the proof.
\end{proof}

\section{Applications to cryptography}

\begin{de}
Cryptography studies methods for secure communication in the presence of an adversary.  A plaintext $M$ is converted by an encryption map $E$ to a ciphertext $C=E(M)$; a decryption map $D$ should satisfy $D(C)=M$ for the intended recipient.
\end{de}

\begin{app}[Caesar cipher]
Encode letters as $0,1,\ldots,25$.  For a key $k$, encrypt and decrypt componentwise by
\[
 b_i\equiv a_i+k\pmod{26},
 \qquad
 a_i\equiv b_i-k\pmod{26}.
\]
This symmetric-key cipher is insecure because the key space is tiny and letter frequencies are preserved.
\end{app}

\begin{app}[Diffie--Hellman key exchange]
Publicly choose a large prime $p$ and a primitive root $g$ modulo $p$.  Alice chooses secret $a$, publishes $A=g^a\bmod p$; Bob chooses secret $b$, publishes $B=g^b\bmod p$.  They independently compute
\[
 B^a\equiv A^b\equiv g^{ab}\pmod p.
\]
Security rests on the practical difficulty of the discrete logarithm problem for suitable parameters.
\end{app}

\begin{app}[RSA]
Choose distinct large primes $p,q$, set $n=pq$, and choose $e$ with
$\gcd(e,\totient(n))=1$.  Let $d$ satisfy
\[
 ed\equiv1\pmod{\totient(n)}.
\]
The public key is $(n,e)$ and the private exponent is $d$.  For a message representative $0\le M<n$,
\[
 C\equiv M^e\pmod n,
 \qquad
 M\equiv C^d\pmod n.
\]
Indeed, write $ed=1+k\totient(n)$.  Modulo $p$, either $p\mid M$, in which case both $M^{ed}$ and $M$ are $0$, or Fermat's theorem gives
$M^{ed}=M(M^{p-1})^{k(q-1)}\equiv M$.  The same argument holds modulo $q$, so the Chinese remainder theorem gives $M^{ed}\equiv M\pmod n$ for every $M$, without requiring $\gcd(M,n)=1$.

For real systems, messages are encoded with randomized padding rather than textbook exponentiation.  For a semiprime $n=pq$, knowing $\totient(n)$ reveals
$p+q=n+1-\totient(n)$ and hence factors $n$ by solving a quadratic; factoring also immediately gives $\totient(n)$.
\end{app}

\chapter{Quadratic Residues}

\begin{notedefinition}[5.1]
Let $p$ be an odd prime and let $\gcd(a,p)=1$.  The integer $a$ is a quadratic residue modulo $p$ if
\[
 x^2\equiv a\pmod p
\]
has a solution, and a quadratic nonresidue otherwise.
\end{notedefinition}
\begin{noteremark}
A nonzero square has exactly two square roots modulo $p$, namely $\pm x_0$.  Hence there are $(p-1)/2$ nonzero quadratic residues and $(p-1)/2$ nonresidues.  If $g$ is a primitive root modulo $p$, the residues are
$g^2,g^4,\ldots,g^{p-1}$.
\end{noteremark}

\begin{notetheorem}[5.1 (Euler's criterion)]
Let $p$ be an odd prime and $\gcd(a,p)=1$.  Then $a$ is a quadratic residue modulo $p$ if and only if
\[
 a^{(p-1)/2}\equiv1\pmod p.
\]
\end{notetheorem}
\begin{proof}
Apply Theorem 4.19 with $m=p$ and $k=2$.
\end{proof}

\begin{notecorollary}[5.2]
Under the same hypotheses, $a$ is a quadratic nonresidue modulo $p$ if and only if
\[
 a^{(p-1)/2}\equiv-1\pmod p.
\]
\end{notecorollary}

\begin{notedefinition}[5.2]
For an odd prime $p$, the Legendre symbol is
\[
 \legendre ap=
 \begin{cases}
 0,&p\mid a,\\
 1,&a\text{ is a quadratic residue modulo }p,\\
 -1,&a\text{ is a quadratic nonresidue modulo }p.
 \end{cases}
\]
\end{notedefinition}

\begin{notetheorem}[5.2]
For an odd prime $p$ and integers $a,b$:
\begin{enumerate}[label=(\alph*)]
\item $\legendre ap\equiv a^{(p-1)/2}\pmod p$;
\item if $a\equiv b\pmod p$, then $\legendre ap=\legendre bp$;
\item $\legendre{ab}p=\legendre ap\legendre bp$;
\item if $p\nmid a$, then $\legendre{a^2}p=1$.
\end{enumerate}
\end{notetheorem}
\begin{proof}
Part (a) is Euler's criterion, including the case $p\mid a$.  Part (b) is immediate.  Multiplying the congruences in (a) proves (c), because both sides lie in $\{-1,0,1\}$; part (d) follows directly.
\end{proof}

\begin{noteremark}
Euler's criterion gives
\[
 \legendre{-1}p=
 \begin{cases}
 1,&p\equiv1\pmod4,\\
 -1,&p\equiv3\pmod4.
 \end{cases}
\]
\end{noteremark}

\begin{notetheorem}[5.3 (Gauss's lemma)]
Let $p$ be an odd prime and $p\nmid a$.  For
\[
 a,2a,\ldots,\frac{p-1}{2}a,
\]
let $n$ be the number of least positive residues that exceed $p/2$.  Then
\[
 \legendre ap=(-1)^n.
\]
\end{notetheorem}
\begin{proof}
Replace every least positive residue exceeding $p/2$ by its negative.  The resulting absolute residues are a permutation of
$1,2,\ldots,(p-1)/2$: no two can agree up to sign unless the original indices agree.  Hence
\[
 a^{(p-1)/2}\left(\frac{p-1}{2}\right)!
 \equiv(-1)^n\left(\frac{p-1}{2}\right)!\pmod p.
\]
Cancel the factorial and apply Euler's criterion.
\end{proof}

\begin{notetheorem}[5.4]
Let $p$ be an odd prime and $p\nmid a$.
\begin{enumerate}[label=(\alph*)]
\item If $a$ is odd, then
\[
 \legendre ap=(-1)^{\sum_{k=1}^{(p-1)/2}\floor{ka/p}}.
\]
\item
\[
 \legendre2p=(-1)^{(p^2-1)/8}.
\]
\end{enumerate}
\end{notetheorem}
\begin{proof}
For (a), write $ka=pq_k+r_k$ with $1\le r_k\le p-1$.  Let $n$ be the number of $r_k>p/2$, and replace those $r_k$ by $p-r_k$.  These absolute residues are a permutation of $1,\ldots,(p-1)/2$.  Comparing
\[
 a\sum k=p\sum q_k+\sum r_k
\]
modulo $2$, and using that $a$ and $p$ are odd, gives
$n\equiv\sum q_k\pmod2$.  Gauss's lemma proves (a).

For (b), Gauss's lemma counts the integers $k$ with
$1\le k\le(p-1)/2$ and $2k>p/2$.  Their number is
\[
 \frac{p-1}{2}-\floor{p/4},
\]
whose parity is $(p^2-1)/8$; this may be checked in the four cases $p\equiv1,3,5,7\pmod8$.
\end{proof}

\begin{notecorollary}[5.5 (second supplementary law)]
For an odd prime $p$,
\[
 \legendre2p=
 \begin{cases}
 1,&p\equiv1,7\pmod8,\\
 -1,&p\equiv3,5\pmod8.
 \end{cases}
\]
\end{notecorollary}

\begin{notetheorem}[5.6]
Let $p$ and $q=2p+1$ be odd primes.  Then
\[
 2(-1)^{(p-1)/2}
\]
is a primitive root modulo $q$.
\end{notetheorem}
\begin{proof}
Let $x=2(-1)^{(p-1)/2}$ and let $m=\ord_q(x)$.  Since $m\mid2p$ and $x\not\equiv\pm1\pmod q$, the only possible nonprimitive order is $m=p$.

If $p\equiv1\pmod4$, then $q\equiv3\pmod8$ and $x=2$.  Thus
\[
 x^p=2^{(q-1)/2}\equiv\legendre2q=-1\pmod q,
\]
contradicting $m=p$.  If $p\equiv3\pmod4$, then $q\equiv7\pmod8$ and $x=-2$; now
\[
 x^p=-2^{(q-1)/2}\equiv-\legendre2q=-1\pmod q,
\]
again a contradiction.  Hence $m=2p=q-1$.
\end{proof}

\begin{notetheorem}[5.7 (quadratic reciprocity)]
Let $p$ and $q$ be distinct odd primes.  Then
\[
 \legendre pq\legendre qp
 =(-1)^{\frac{p-1}{2}\frac{q-1}{2}}.
\]
\end{notetheorem}
\begin{proof}
By Theorem 5.4(a),
\[
 \legendre pq=(-1)^{\sum_{k=1}^{(q-1)/2}\floor{kp/q}},
 \qquad
 \legendre qp=(-1)^{\sum_{j=1}^{(p-1)/2}\floor{jq/p}}.
\]
In the rectangle
$1\le x\le(q-1)/2$, $1\le y\le(p-1)/2$, no lattice point lies on the line $qy=px$, because $p$ and $q$ are distinct primes.  The first floor sum counts lattice points on one side of the line and the second counts those on the other.  Their sum is therefore
$((p-1)/2)((q-1)/2)$, proving the identity.
\end{proof}

\begin{notecorollary}[5.8]
For distinct odd primes $p,q$,
\[
 \legendre pq=
 \begin{cases}
 -\legendre qp,&p\equiv q\equiv3\pmod4,\\
 \phantom{-}\legendre qp,&\text{otherwise}.
 \end{cases}
\]
\end{notecorollary}

\begin{notetheorem}[5.9]
Let $p$ be an odd prime, $k\ge1$, and $\gcd(a,p)=1$.  The congruence
\[
 x^2\equiv a\pmod{p^k}
\]
has a solution if and only if $\legendre ap=1$.
\end{notetheorem}
\begin{proof}
A solution modulo $p^k$ reduces to a solution modulo $p$, proving necessity.  Conversely, a root $x_1$ modulo $p$ satisfies
$f'(x_1)=2x_1\not\equiv0\pmod p$ for $f(x)=x^2-a$.  Hensel's lemma lifts it uniquely through every power of $p$.
\end{proof}

\begin{notetheorem}[5.10]
Let $a$ be odd and $k\ge3$.  The congruence
\[
 x^2\equiv a\pmod{2^k}
\]
has a solution if and only if
\[
 a\equiv1\pmod8.
\]
\end{notetheorem}
\begin{proof}
Every odd square is $1$ modulo $8$, proving necessity.  For sufficiency, induct on $k$.  The case $k=3$ is immediate.  Suppose $x_0^2\equiv a\pmod{2^k}$ with $x_0$ odd.  The two candidates
\[
 x_0\quad\text{and}\quad x_0+2^{k-1}
\]
have squares congruent modulo $2^k$ and differ modulo $2^{k+1}$ by
$x_0 2^k\equiv2^k\pmod{2^{k+1}}$.  Exactly one therefore has square congruent to $a$ modulo $2^{k+1}$.
\end{proof}

\section{Jacobi symbols}

\begin{notedefinition}[5.3]
Let $n$ be a positive odd integer with prime factorization
$n=p_1^{\alpha_1}\cdots p_r^{\alpha_r}$.  The Jacobi symbol is
\[
 \jacobi an=\prod_{i=1}^r\legendre a{p_i}^{\alpha_i},
 \qquad
 \jacobi a1=1.
\]
\end{notedefinition}
\begin{noteremark}
For prime denominator the Jacobi symbol is the Legendre symbol.  For composite denominator, the value $1$ does not guarantee that $a$ is a square modulo $n$; for example, $\jacobi2{15}=1$, but $x^2\equiv2\pmod{15}$ has no solution.
\end{noteremark}

\begin{notetheorem}[5.11]
Let $n$ be odd and $\gcd(a,n)=1$.  If $x^2\equiv a\pmod n$ is solvable, then
\[
 \jacobi an=1.
\]
\end{notetheorem}
\begin{proof}
A solution modulo $n$ is a solution modulo every prime divisor $p_i$, so each Legendre symbol is $1$.
\end{proof}

\begin{notetheorem}[5.12]
For positive odd $n$ and integers $a,b$:
\begin{enumerate}[label=(\alph*)]
\item if $a\equiv b\pmod n$, then $\jacobi an=\jacobi bn$;
\item $\jacobi{ab}n=\jacobi an\jacobi bn$;
\item $\jacobi{-1}n=(-1)^{(n-1)/2}$;
\item $\jacobi2n=(-1)^{(n^2-1)/8}$.
\end{enumerate}
\end{notetheorem}
\begin{proof}
Parts (a) and (b) follow prime factor by prime factor.  Parts (c) and (d) follow by multiplying the corresponding supplementary laws and using the parity identities in Lemma 5.13.
\end{proof}

\begin{notelemma}[5.13]
For odd integers $a,b$,
\begin{align*}
 \frac{a-1}{2}+\frac{b-1}{2}
 &\equiv\frac{ab-1}{2}\pmod2,\\
 \frac{a^2-1}{8}+\frac{b^2-1}{8}
 &\equiv\frac{a^2b^2-1}{8}\pmod2.
\end{align*}
\end{notelemma}
\begin{proof}
The differences are $(a-1)(b-1)/2$ and
$(a^2-1)(b^2-1)/8$, respectively, and both are even for odd $a,b$.
\end{proof}

\begin{notetheorem}[5.14 (quadratic reciprocity for Jacobi symbols)]
If $m,n$ are coprime positive odd integers, then
\[
 \jacobi mn\jacobi nm
 =(-1)^{\frac{m-1}{2}\frac{n-1}{2}}.
\]
Equivalently, the symbols change sign exactly when $m\equiv n\equiv3\pmod4$.
\end{notetheorem}
\begin{proof}
Factor $m$ and $n$ into primes, apply prime quadratic reciprocity to every pair, and collect the signs.  Lemma 5.13 combines the resulting exponents into the displayed one.
\end{proof}

\begin{noteremark}[(computing symbols)]
Repeated use of numerator reduction, multiplicativity, the two supplementary laws, and reciprocity computes Jacobi and Legendre symbols without factoring the numerator.  Euler's criterion, however, does not extend from Legendre to Jacobi symbols: for example,
$2^{170}\equiv1\pmod{341}$ while $\jacobi2{341}=-1$.
\end{noteremark}

\section{The Solovay--Strassen primality test}

\begin{noteapplication}[5.1]
For odd $n>2$, choose a random $b\in\{2,\ldots,n-1\}$.  If
$\gcd(b,n)>1$, return ``composite.''  Otherwise compute the Jacobi symbol and test
\[
 b^{(n-1)/2}\equiv\jacobi bn\pmod n.
\]
Failure proves that $n$ is composite; repeated success gives ``probably prime.''
\end{noteapplication}

\begin{notelemma}[5.16]
If $n$ is a positive odd integer that is not a perfect square, then there exists a unit $b$ modulo $n$ such that
\[
 \jacobi bn=-1.
\]
\end{notelemma}
\begin{proof}
Write $n=\prod p_i^{\alpha_i}$ and choose an index $j$ for which $\alpha_j$ is odd.  Choose a quadratic nonresidue $b_0$ modulo $p_j$, and use the Chinese remainder theorem to choose $b$ satisfying
\[
 b\equiv b_0\pmod{p_j},
 \qquad
 b\equiv1\pmod{p_i}\quad(i\ne j).
\]
Then $b$ is a unit and its Jacobi symbol is $(-1)^{\alpha_j}=-1$.
\end{proof}

\begin{notelemma}[5.17]
If $n$ is odd and composite, then some unit $b$ satisfies
\[
 b^{(n-1)/2}\not\equiv\jacobi bn\pmod n.
\]
\end{notelemma}
\begin{proof}
Assume the congruence holds for every unit.  Squaring it gives
$b^{n-1}\equiv1\pmod n$ for every unit, so $n$ is a Carmichael number.

First, $n$ must be squarefree.  If $p^2\mid n$, choose by the Chinese remainder theorem a unit $b$ whose reduction modulo $p^2$ is a primitive root and whose reductions at the other prime-power factors are $1$.  Then
$b^{n-1}\equiv1\pmod{p^2}$ forces $p(p-1)\mid n-1$, hence $p\mid n-1$, contradicting $p\mid n$.

Thus $n=q_1\cdots q_r$ with $r\ge2$.  For each $q_i$, use the Chinese remainder theorem to choose a unit that is primitive modulo $q_i$ and congruent to $1$ modulo the other prime factors.  The Carmichael congruence then implies
$q_i-1\mid n-1$.

By Lemma 5.16 choose $c$ with $\jacobi cn=-1$.  Our assumption gives
$c^{(n-1)/2}\equiv-1\pmod n$.  Consequently, for every $q_i$, the integer quotient $(n-1)/(q_i-1)$ must be odd; otherwise Fermat's theorem would give $c^{(n-1)/2}\equiv1\pmod{q_i}$.

Choose a primitive root $g$ modulo $q_1$ and use the Chinese remainder theorem to choose $b$ with
$b\equiv g\pmod{q_1}$ and $b\equiv1\pmod{q_i}$ for $i>1$.  A primitive root is a quadratic nonresidue, so $\jacobi bn=-1$.  Because $(n-1)/(q_1-1)$ is odd,
$b^{(n-1)/2}\equiv-1\pmod{q_1}$, while it is $1$ modulo every other $q_i$.  Hence it is not congruent to $-1$ modulo $n$, contradicting the assumed Euler congruence.
\end{proof}

\begin{notetheorem}[5.15]
Let $n$ be odd and composite.  Among the $\totient(n)$ units modulo $n$, at most $\totient(n)/2$ are Euler liars for the Solovay--Strassen test.
\end{notetheorem}
\begin{proof}
On the unit group define
\[
 \Psi(b)=\jacobi bn\, b^{-(n-1)/2}\pmod n.
\]
Multiplicativity of the Jacobi symbol makes $\Psi$ a group homomorphism.  Its kernel is exactly the set of Euler liars.  Lemma 5.17 says that $\Psi$ is nontrivial, so its kernel has index at least $2$ and therefore size at most $\totient(n)/2$.
\end{proof}

\chapter{Continued Fractions}

\begin{de}
A finite continued fraction is an expression
\[
 [a_0;a_1,\ldots,a_n]
 =a_0+\cfrac1{a_1+\cfrac1{\ddots+\cfrac1{a_n}}},
\]
where $a_i\in\R$ and $a_i>0$ for $i\ge1$.  It is simple if every $a_i$ is an integer.
\end{de}

\begin{thm}
Every finite simple continued fraction is rational.
\end{thm}
\begin{proof}
Evaluate from the final entry upward; rationality is preserved by addition and reciprocal of a nonzero rational.
\end{proof}

\begin{thm}
Every rational number $a/b$, with $a\in\Z$ and $b\in\N$, has a finite simple continued-fraction representation.  Its entries are the quotients in the Euclidean algorithm for $a$ and $b$.
\end{thm}
\begin{proof}
Write $a=a_0b+r$ with $0\le r<b$.  If $r=0$, then $a/b=[a_0]$.  Otherwise
\[
 \frac ab=a_0+\frac1{b/r},
\]
and the denominator in $b/r$ is $r<b$.  Induction on $b$ completes the construction.
\end{proof}

\begin{remark}[the two finite representations]
For an expansion of positive length with $a_n>1$,
\[
 [a_0;\ldots,a_n]=[a_0;\ldots,a_n-1,1].
\]
Conversely, a terminal $1$ may be absorbed into the preceding entry.  Thus every nonintegral rational has exactly two finite simple representations, exactly one of which ends in an entry greater than $1$.  An integer $a_0$ has the two representations $[a_0]$ and $[a_0-1;1]$; by convention the one-term expansion is taken as canonical.
\end{remark}

\begin{thm}
Suppose
\[
 [a_0;\ldots,a_n]=[b_0;\ldots,b_m],
 \qquad a_n>1,\quad b_m>1.
\]
Then $m=n$ and $a_i=b_i$ for every $i$.
\end{thm}
\begin{proof}
For a finite simple continued fraction whose final entry exceeds $1$, every complete quotient after the first is greater than $1$.  Hence the integer part of the value is $a_0$, and the reciprocal of the fractional part is the next complete quotient.  Equality of the two values gives $a_0=b_0$ and equality of the tails.  Repeating proves equality of all entries and simultaneous termination.
\end{proof}

\begin{de}
Let $a_0\in\R$ and $a_i>0$ for $i\ge1$.  The $k$th convergent of the sequence $(a_i)$ is
\[
 C_k=[a_0;a_1,\ldots,a_k].
\]
The infinite continued fraction $[a_0;a_1,a_2,\ldots]$ is the limit of $C_k$ when that limit exists.  It is simple if all $a_i\in\Z$.
\end{de}

\begin{thm}
The convergents satisfy
\[
 C_k=\frac{p_k}{q_k},
\]
where
\[
 p_{-2}=0,\quad p_{-1}=1,\quad p_k=a_kp_{k-1}+p_{k-2},
\]
\[
 q_{-2}=1,\quad q_{-1}=0,\quad q_k=a_kq_{k-1}+q_{k-2}.
\]
\end{thm}
\begin{proof}
Induct on $k$, using
\[
 [a_0;\ldots,a_k]
 =[a_0;\ldots,a_{k-1}+1/a_k]
\]
and the recurrence at the final step.  Equivalently, multiply the matrices
\[
 \begin{pmatrix}a_i&1\\1&0\end{pmatrix}.
\]
\end{proof}

\begin{thm}
For $k\ge1$,
\[
 p_kq_{k-1}-p_{k-1}q_k=(-1)^{k-1},
\]
and for $k\ge2$,
\[
 p_kq_{k-2}-p_{k-2}q_k=(-1)^k a_k.
\]
\end{thm}
\begin{proof}
The first identity follows by induction from the recurrences:
\[
 p_kq_{k-1}-p_{k-1}q_k
 =-(p_{k-1}q_{k-2}-p_{k-2}q_{k-1}).
\]
The second is
\[
 a_k(p_{k-1}q_{k-2}-p_{k-2}q_{k-1}).
\]
\end{proof}

\begin{thm}
For the convergents $C_k=p_k/q_k$,
\[
 C_k-C_{k-1}=\frac{(-1)^{k-1}}{q_kq_{k-1}}
 \quad(k\ge1),
\]
and
\[
 C_k-C_{k-2}=\frac{a_k(-1)^k}{q_kq_{k-2}}
 \quad(k\ge2).
\]
\end{thm}
\begin{proof}
Divide the identities in Theorem 6.5 by the appropriate products of denominators.
\end{proof}

\begin{thm}
For every convergent of a simple continued fraction,
\[
 \gcd(p_k,q_k)=1.
\]
\end{thm}
\begin{proof}
Any common divisor divides
$p_kq_{k-1}-p_{k-1}q_k=\pm1$.
\end{proof}

\begin{thm}
The convergents of an infinite simple continued fraction satisfy
\begin{align*}
 C_1&>C_3>C_5>\cdots,\\
 C_0&<C_2<C_4<\cdots,\\
 C_{2i+1}&>C_{2j}\qquad(i,j\ge0).
\end{align*}
Consequently the limit $\lim_{k\to\infty}C_k$ exists.
\end{thm}
\begin{proof}
The signs in Theorem 6.6 show that the even subsequence is increasing and the odd subsequence decreasing.  They also show each odd convergent exceeds the neighboring even convergents, hence all of them.  The two subsequences are monotone and bounded, and
\[
 C_{2k+1}-C_{2k}=\frac1{q_{2k+1}q_{2k}}\longrightarrow0
\]
because the positive integers $q_k$ increase without bound.  Their limits are therefore equal.
\end{proof}

\begin{thm}
Every infinite simple continued fraction is irrational.
\end{thm}
\begin{proof}
Let its value be $x$.  Suppose $x=p/q$ in lowest terms.  The limit lies strictly between consecutive convergents, so
\[
 0<\abs*{x-C_k}<\abs*{C_{k+1}-C_k}
 =\frac1{q_kq_{k+1}}.
\]
On the other hand,
\[
 \abs*{x-C_k}
 =\frac{\abs*{pq_k-qp_k}}{q q_k}
 \ge\frac1{q q_k}.
\]
Hence $q>q_{k+1}$ for every $k$, impossible.
\end{proof}

\begin{thm}
Two infinite simple continued fractions are equal if and only if all corresponding entries are equal.
\end{thm}
\begin{proof}
If $x=[a_0;a_1,\ldots]$, then $a_0=\floor x$ because the tail after $a_0$ lies in $(0,1)$.  The next complete quotient is
$1/(x-a_0)$.  Equality of two values therefore determines the same first entry and equal tails; induction gives every entry.
\end{proof}

\begin{thm}
Every irrational real number $x$ has an infinite simple continued-fraction expansion.
\end{thm}
\begin{proof}
Set $x_0=x$ and define recursively
\[
 a_k=\floor{x_k},
 \qquad
 x_{k+1}=\frac1{x_k-a_k}.
\]
Every $x_k$ is irrational, so the denominator is nonzero; for $k\ge1$, $x_k>1$, hence $a_k\ge1$.

The finite-tail identity is
\[
 x=\frac{p_nx_{n+1}+p_{n-1}}{q_nx_{n+1}+q_{n-1}}.
\]
Using Theorem 6.5,
\[
 \abs*{x-\frac{p_n}{q_n}}
 =\frac1{q_n(q_nx_{n+1}+q_{n-1})}
 <\frac1{q_n^2}\longrightarrow0.
\]
Thus the convergents tend to $x$.
\end{proof}

\section{Best approximation}

\begin{thm}
Let $x$ be irrational and let $p_k/q_k$ be its $k$th convergent.  If $a,b\in\Z$ and
\[
 1\le b<q_{k+1},
\]
then
\[
 \abs*{q_kx-p_k}\le\abs*{bx-a}.
\]
\end{thm}
\begin{proof}
Because
$p_kq_{k+1}-p_{k+1}q_k=\pm1$, there are unique integers $r,s$ such that
\[
 a=rp_k+sp_{k+1},
 \qquad
 b=rq_k+sq_{k+1}.
\]
If $s=0$, the conclusion is immediate.  The case $r=0$ is impossible because then $b\ge q_{k+1}$.  If $r$ and $s$ had the same sign, positivity of $b$ would give $b\ge q_k+q_{k+1}$.  Thus they have opposite signs.

The quantities $q_kx-p_k$ and $q_{k+1}x-p_{k+1}$ also have opposite signs.  Therefore the two terms in
\[
 bx-a=r(q_kx-p_k)+s(q_{k+1}x-p_{k+1})
\]
have the same sign, and
\[
 \abs*{bx-a}
 =\abs r\,\abs*{q_kx-p_k}+\abs s\,\abs*{q_{k+1}x-p_{k+1}}
 \ge\abs*{q_kx-p_k}.
\]
\end{proof}

\begin{thm}
Under the same hypotheses, if $a/b\in\Q$ and $1\le b\le q_k$, then
\[
 \abs*{x-\frac{p_k}{q_k}}
 \le\abs*{x-\frac ab}.
\]
\end{thm}
\begin{proof}
Since $b<q_{k+1}$, Theorem 6.12 gives
\[
 \abs*{x-p_k/q_k}
 =\frac{\abs*{q_kx-p_k}}{q_k}
 \le\frac{\abs*{bx-a}}{q_k}
 \le\frac{\abs*{bx-a}}b.
\]
\end{proof}

\begin{thm}[Legendre's criterion]
Let $x$ be irrational and let $a/b\in\Q$ be in lowest terms with $b\ge1$.  If
\[
 \abs*{x-\frac ab}<\frac1{2b^2},
\]
then $a/b$ is a convergent of $x$.
\end{thm}
\begin{proof}
Suppose not.  Choose $k$ with $q_k\le b<q_{k+1}$.  Theorem 6.12 gives
\[
 \abs*{x-\frac{p_k}{q_k}}
 \le\frac{\abs*{bx-a}}{q_k}
 <\frac1{2bq_k}.
\]
Since $a/b$ and $p_k/q_k$ are distinct reduced fractions,
\[
 \frac1{bq_k}
 \le\abs*{\frac ab-\frac{p_k}{q_k}}
 \le\abs*{x-\frac ab}+\abs*{x-\frac{p_k}{q_k}}
 <\frac1{2b^2}+\frac1{2bq_k}
 \le\frac1{bq_k},
\]
a contradiction.
\end{proof}

\section{Periodic continued fractions and quadratic irrationals}

\begin{de}
An infinite continued fraction $[a_0;a_1,a_2,\ldots]$ is periodic if there exist $N\ge0$ and $r\ge1$ such that $a_{n+r}=a_n$ for all $n\ge N$.  We write
\[
 [a_0;\ldots,a_{N-1},\overbar{a_N,\ldots,a_{N+r-1}}].
\]
\end{de}

\begin{de}
An irrational real number $\alpha$ is a quadratic irrational if
\[
 A\alpha^2+B\alpha+C=0
\]
for integers $A,B,C$ with $A\ne0$.
\end{de}

\begin{thm}
A real number $\alpha$ is a quadratic irrational if and only if
\[
 \alpha=\frac{a+\sqrt d}{c}
\]
for integers $a,c$ and a positive nonsquare integer $d$, with $c\ne0$.
\end{thm}
\begin{proof}
If $A\alpha^2+B\alpha+C=0$, then
\[
 \alpha=\frac{-B\pm\sqrt{B^2-4AC}}{2A}.
\]
Because $\alpha$ is real and irrational, the discriminant is positive and not a perfect square; the sign may be absorbed into $a,c$.

Conversely, $(c\alpha-a)^2=d$, so
\[
 c^2\alpha^2-2ac\alpha+(a^2-d)=0.
\]
The number is irrational because $d$ is not a square.
\end{proof}

\begin{thm}
Let $r,s,t,u\in\Z$ with $t$ and $u$ not both zero, and let $\alpha$ be a quadratic irrational.  Then
\[
 \frac{r\alpha+s}{t\alpha+u}
\]
is rational or a quadratic irrational.
\end{thm}
\begin{proof}
The denominator is nonzero, since otherwise $\alpha$ would be rational.  Write $\alpha=(a+\sqrt d)/c$ and rationalize.  The result has the form $A+B\sqrt d$ with $A,B\in\Q$.  If $B=0$ it is rational; otherwise it satisfies a quadratic equation over $\Q$, which may be cleared of denominators.
\end{proof}

\begin{thm}
Every periodic simple continued fraction represents a quadratic irrational.
\end{thm}
\begin{proof}
Let the repeating tail be
\[
 \alpha=[\overbar{a_N,\ldots,a_{N+r-1}}].
\]
The finite-tail formula expresses one period as a linear fractional transformation:
\[
 \alpha=\frac{p_{r-1}\alpha+p_{r-2}}{q_{r-1}\alpha+q_{r-2}}.
\]
After clearing the denominator, $\alpha$ satisfies a quadratic equation with integer coefficients.  It is irrational by Theorem 6.9, hence quadratic irrational.  The original value is a linear fractional transformation of $\alpha$, so Theorem 6.16 and Theorem 6.9 show that it too is a quadratic irrational.
\end{proof}

\begin{de}
If
\[
 \alpha=\frac{a+c\sqrt d}{b}
\]
is quadratic irrational, its conjugate is
\[
 \alpha^*=\frac{a-c\sqrt d}{b}.
\]
Conjugation commutes with addition, multiplication, and division by a nonzero quadratic irrational.
\end{de}

\begin{thm}[Lagrange]
Every quadratic irrational has a periodic simple continued-fraction expansion.
\end{thm}
\begin{proof}
After multiplying numerator, denominator, and radicand by a suitable square, write the initial complete quotient in the standard form
\[
 x_0=\frac{P_0+\sqrt D}{Q_0},
\]
where $D$ is a positive nonsquare integer, $P_0,Q_0\in\Z$, and
$Q_0\mid D-P_0^2$.  If
$x_k=(P_k+\sqrt D)/Q_k$ and $a_k=\floor{x_k}$, then
\[
 P_{k+1}=a_kQ_k-P_k,
 \qquad
 Q_{k+1}=\frac{D-P_{k+1}^2}{Q_k},
\]
so induction gives $P_k,Q_k\in\Z$ and $Q_k\ne0$.

Let $p_j/q_j$ be the convergents of $x_0$.  Solving the finite-tail identity for $x_k$ and conjugating gives
\[
 x_k^*=\frac{p_{k-2}-x_0^*q_{k-2}}{x_0^*q_{k-1}-p_{k-1}}.
\]
Since $p_j/q_j\to x_0$ and $x_0^*\ne x_0$, the numerator and denominator eventually have opposite signs.  Hence $x_k^*<0$ for all sufficiently large $k$.  Since $x_k>1$, it follows that eventually
\[
 Q_k>0,
 \qquad
 -\sqrt D<P_k<\sqrt D.
\]
Moreover, $Q_k$ is a positive divisor of $D-P_k^2$, so $Q_k\le D$.  Only finitely many integer pairs $(P_k,Q_k)$ can occur.  Some pair therefore repeats, and the recurrence then forces all subsequent complete quotients and partial quotients to repeat.  Thus the continued fraction is periodic.
\end{proof}

\begin{de}
A continued fraction is purely periodic if it has the form
$[\overbar{a_0,\ldots,a_{r-1}}]$.  A quadratic irrational $x$ is reduced if
\[
 x>1,
 \qquad
 -1<x^*<0.
\]
\end{de}

\begin{thm}
A real number is a reduced quadratic irrational if and only if its simple continued-fraction expansion is purely periodic.
\end{thm}
\begin{proof}
Suppose first that $x$ is reduced.  If $x_k$ is a reduced complete quotient and $a_k=\floor{x_k}$, then
\[
 x_{k+1}=\frac1{x_k-a_k}>1,
 \qquad
 x_{k+1}^*=\frac1{x_k^*-a_k}\in(-1,0).
\]
Thus every complete quotient is reduced.  Lagrange's theorem gives $x_N=x_{N+r}$ for some $N,r$.  A reduced quotient $y$ has a unique reduced predecessor: from
$z=a+1/y$ and $z^*=a+1/y^*\in(-1,0)$ one obtains
\[
 a=\floor{-1/y^*}.
\]
Therefore equality of two reduced complete quotients propagates backward, giving $x_0=x_r$.  The expansion is purely periodic.

Conversely, let
$x=[\overbar{a_0,\ldots,a_{r-1}}]$.  Then $a_0\ge1$ and $x>1$.  The period relation gives
\[
 q_{r-1}x^2+(q_{r-2}-p_{r-1})x-p_{r-2}=0.
\]
Let $f$ be the left side.  We have $f(0)=-p_{r-2}<0$ and
\[
 f(-1)=(q_{r-1}-q_{r-2})+(p_{r-1}-p_{r-2})>0.
\]
Thus the other root $x^*$ lies in $(-1,0)$, and $x$ is reduced.
\end{proof}

\begin{remark}[square roots]
If $D$ is a positive nonsquare and $a_0=\floor{\sqrt D}$, then
\[
 \sqrt D=[a_0;\overbar{a_1,\ldots,a_{r-1},2a_0}].
\]
This follows by applying pure periodicity to $a_0+\sqrt D$.
\end{remark}

\begin{thm}
Let $D$ be a positive nonsquare integer, let $p_k/q_k$ be the $k$th convergent of $\sqrt D$, and write the complete quotients as
\[
 x_{k+1}=\frac{P_{k+1}+\sqrt D}{Q_{k+1}}.
\]
Then for every $k\ge0$,
\[
 p_k^2-Dq_k^2=(-1)^{k+1}Q_{k+1}.
\]
\end{thm}
\begin{proof}
The finite-tail identity
\[
 \sqrt D=\frac{p_kx_{k+1}+p_{k-1}}{q_kx_{k+1}+q_{k-1}}
\]
and comparison of rational and irrational parts give
\[
 p_k=q_kP_{k+1}+q_{k-1}Q_{k+1},
 \qquad
 Dq_k=p_kP_{k+1}+p_{k-1}Q_{k+1}.
\]
Therefore
\begin{align*}
 p_k^2-Dq_k^2
 &=Q_{k+1}(p_kq_{k-1}-p_{k-1}q_k)\\
 &=(-1)^{k-1}Q_{k+1}=(-1)^{k+1}Q_{k+1}.
\end{align*}
\end{proof}

\chapter{Nonlinear Diophantine Equations}

\begin{de}
Let $d\in\N$ be nonsquare.  The equations
\[
 x^2-dy^2=1
 \qquad\text{and}\qquad
 x^2-dy^2=-1
\]
are the positive and negative Pell equations.
\end{de}

\begin{lem}
Let $x_0,y_0\in\N$, let $\delta>0$ with $\sqrt\delta\notin\Q$, and suppose
\[
 x_0^2-\delta y_0^2=\alpha,
 \qquad
 0\le\alpha<\sqrt\delta.
\]
Then
\[
 \abs*{\frac{x_0}{y_0}-\sqrt\delta}<\frac1{2y_0^2}.
\]
\end{lem}
\begin{proof}
The case $\alpha=0$ would make $\sqrt\delta=x_0/y_0$ rational, so in fact $\alpha>0$.  Hence $x_0^2>\delta y_0^2$, $x_0>\sqrt\delta y_0$, and
\[
 \frac{x_0}{y_0}-\sqrt\delta
 =\frac{\alpha}{y_0(x_0+\sqrt\delta y_0)}
 <\frac{\sqrt\delta}{2\sqrt\delta y_0^2}.
\]
\end{proof}

\begin{thm}
Let $d$ be a positive nonsquare integer and let $n\in\Z$ satisfy $\abs n<\sqrt d$.  If $(x_0,y_0)\in\N^2$ satisfies
\[
 x_0^2-dy_0^2=n,
\]
then $x_0/y_0$ is a convergent of $\sqrt d$.
\end{thm}
\begin{proof}
If $n>0$, Lemma 7.1 and Legendre's criterion apply directly.  If $n<0$, rewrite the equation as
\[
 y_0^2-\frac1d x_0^2=-\frac nd,
 \qquad
 0<-\frac nd<\frac1{\sqrt d}.
\]
After reducing the fraction if necessary, Lemma 7.1 and Legendre's criterion show that $y_0/x_0$ is a convergent of $1/\sqrt d$.  The continued fraction of $1/x$ is obtained by adding a leading zero, so its nonzero convergents are the reciprocals of the convergents of $x$.  Hence $x_0/y_0$ is a convergent of $\sqrt d$.  The case $n=0$ is impossible because $d$ is nonsquare.
\end{proof}

For the rest of the Pell discussion, put
\[
 x_0=\floor{\sqrt d}+\sqrt d
\]
and let its purely periodic continued fraction have least period $r$.  Let
$(P_i,Q_i)$ be the complete-quotient data of Chapter 6, with $Q_0=1$.

\begin{thm}
\begin{enumerate}[label=(\alph*)]
\item $Q_i=1$ if and only if $i\equiv0\pmod r$;
\item $Q_i\ne-1$ for every $i$.
\end{enumerate}
\end{thm}
\begin{proof}
The number $x_0$ is reduced, so its complete quotients are reduced and repeat with least period $r$.  Hence $i\equiv0\pmod r$ implies $Q_i=1$.  Conversely, if $Q_i=1$, then
$x_i=P_i+\sqrt d$.  Reducedness gives
$-1<P_i-\sqrt d<0$, so $P_i=\floor{\sqrt d}$ and $x_i=x_0$.  Minimality of the period gives $r\mid i$.

If $Q_i=-1$, then
$x_i=-P_i-\sqrt d>1$, whereas its conjugate $x_i^*=-P_i+\sqrt d$ is positive, contradicting reducedness.
\end{proof}

\begin{thm}
Let $p_k/q_k$ be the convergents of $\sqrt d$ and let $r$ be its least period.
\begin{enumerate}[label=(\alph*)]
\item If $r$ is even, the negative Pell equation has no positive solution, and the positive solutions are
\[
 (p_{tr-1},q_{tr-1}),\qquad t=1,2,3,\ldots.
\]
\item If $r$ is odd, the negative solutions are the same pairs with $t$ odd, and the positive solutions are those with $t$ even.
\end{enumerate}
\end{thm}
\begin{proof}
By Theorem 7.2, a solution of $x^2-dy^2=\pm1$ must be a convergent.  Theorem 6.20 gives
\[
 p_k^2-dq_k^2=(-1)^{k+1}Q_{k+1}.
\]
Theorem 7.3 shows that this equals $\pm1$ exactly when $k+1=tr$.  The sign is then $(-1)^{tr}$, which gives the two cases.
\end{proof}

\begin{de}
Among the positive solutions to $x^2-dy^2=1$, the one with smallest $x$ is the fundamental solution.
\end{de}

\begin{thm}
Let $(x_1,y_1)$ be the fundamental solution of $x^2-dy^2=1$.  All positive solutions are given by
\[
 x_n+y_n\sqrt d=(x_1+y_1\sqrt d)^n,
 \qquad n\in\N.
\]
\end{thm}
\begin{proof}
The displayed powers have integer coefficients and norm $1$, so they are solutions.  Conversely, let
$\eta=s+t\sqrt d>1$ be any positive solution and put
$\varepsilon=x_1+y_1\sqrt d$.  Choose $n\ge0$ such that
\[
 \varepsilon^n\le\eta<\varepsilon^{n+1}.
\]
Then
\[
 \delta=\eta\varepsilon^{-n}=A+B\sqrt d
\]
has integer coefficients, norm $1$, and $1\le\delta<\varepsilon$.  Its conjugate is $\delta^{-1}>0$, so
$A=(\delta+\delta^{-1})/2>0$ and $B\ge0$.  If $B>0$, then $(A,B)$ is a positive Pell solution.  Since the function $z+z^{-1}$ is increasing for $z\ge1$,
\[
 A=\frac{\delta+\delta^{-1}}2
 <\frac{\varepsilon+\varepsilon^{-1}}2=x_1,
\]
contradicting the definition of the fundamental solution.  Thus $B=0$, whence $\delta=1$ and $\eta=\varepsilon^n$.
\end{proof}

\begin{app}
\begin{enumerate}[label=(\alph*)]
\item There are infinitely many triples of consecutive integers, each of which is a sum of two integer squares.  For every positive solution of $x^2-2y^2=1$,
\[
 x^2-1=y^2+y^2,
 \qquad
 x^2=x^2+0^2,
 \qquad
 x^2+1=x^2+1^2.
\]
\item If $x^2-dy^2=n$ has a nonzero integer solution $(s,t)$ and the positive Pell equation for $d$ has solutions $(x_m,y_m)$, then
\[
 (s+t\sqrt d)(x_m+y_m\sqrt d)=S_m+T_m\sqrt d
\]
gives infinitely many solutions $S_m^2-dT_m^2=n$.
\item There are infinitely many positive integers $n$ for which $n^2+(n+1)^2$ is a square.  Indeed,
\[
 n^2+(n+1)^2=m^2
 \quad\Longleftrightarrow\quad
 (2n+1)^2-2m^2=-1,
\]
and the negative Pell equation for $d=2$ has infinitely many solutions.  Geometrically, these are infinitely many integer right triangles with consecutive legs.
\end{enumerate}
\end{app}

\section{Pythagorean triples}

\begin{de}
A triple $(x,y,z)\in\N^3$ is Pythagorean if $x^2+y^2=z^2$.  It is primitive if $\gcd(x,y,z)=1$.
\end{de}
\begin{remark}
In a primitive Pythagorean triple, $\gcd(x,y)=\gcd(y,z)=\gcd(z,x)=1$, and exactly one of $x,y$ is even.
\end{remark}

\begin{thm}
The primitive Pythagorean triples with $y$ even are precisely
\[
 (x,y,z)=(r^2-s^2,\ 2rs,\ r^2+s^2),
\]
where $r>s$ are coprime positive integers of opposite parity.
\end{thm}
\begin{proof}
Let $(x,y,z)$ be primitive with $y$ even.  Then $x,z$ are odd and
\[
 \left(\frac{z+x}{2}\right)
 \left(\frac{z-x}{2}\right)
 =\left(\frac y2\right)^2.
\]
The two factors on the left are coprime, so each is a square:
$(z+x)/2=r^2$ and $(z-x)/2=s^2$.  This gives the displayed formulas.  Coprimality and opposite parity follow from primitiveness and the oddness of $x$.  Conversely, the formulas satisfy the equation, and the stated conditions make the triple primitive.
\end{proof}

\begin{remark}[rational points on the circle]
The rational points on $X^2+Y^2=1$ are parametrized by lines of rational slope through $(1,0)$:
\[
 (X,Y)=\left(\frac{m^2-1}{m^2+1},\frac{-2m}{m^2+1}\right),
 \qquad m\in\Q.
\]
Clearing denominators recovers the Pythagorean parametrization.
\end{remark}

\section{Infinite descent and local obstructions}

\begin{de}
Fermat's method of infinite descent proves nonexistence by assuming a positive integer solution, constructing a strictly smaller positive solution, and contradicting the well-ordering principle.
\end{de}

\begin{app}
\begin{enumerate}[label=(\alph*)]
\item The system
\[
 x^2+6y^2=z^2,
 \qquad
 6x^2+y^2=w^2
\]
has no positive integer solution.  Otherwise choose one minimizing $x+y+z+w$.  Adding gives
$z^2+w^2=7(x^2+y^2)$.  Since $-1$ is a nonresidue modulo $7$, this congruence forces $7\mid z,w$, and then also $7\mid x,y$.  Dividing all four entries by $7$ produces a smaller positive solution.

\item The equation
\[
 x^3+2y^3=4z^3
\]
has only the zero integer solution.  Modulo $2$, any solution has $x$ even; after substituting $x=2x_1$, one finds $y$ even, and then $z$ even.  Division by $2$ gives a smaller solution of the same equation, so infinite descent forces $x=y=z=0$.
\end{enumerate}
\end{app}

\begin{thm}[Fermat for exponent four]
The equation
\[
 x^4+y^4=z^4
\]
has no nontrivial integer solutions, where nontrivial means $xyz\ne0$.  In fact, the stronger equation
\[
 x^4+y^4=z^2
\]
has no positive integer solutions.
\end{thm}
\begin{proof}
Assume a positive solution to $x^4+y^4=z^2$ with $z$ minimal, and divide out a common factor so that it is primitive.  Exactly one of $x,y$ is even; suppose $y$ is even.  The primitive Pythagorean triple $(x^2,y^2,z)$ gives coprime $r>s$ of opposite parity with
\[
 x^2=r^2-s^2,
 \qquad
 y^2=2rs,
 \qquad
 z=r^2+s^2.
\]
The even parameter must be $s$: if $r$ were even, then $x^2=r^2-s^2\equiv3\pmod4$.  The square condition $2rs=y^2$ and coprimality now imply
\[
 r=m^2,
 \qquad
 s=2n^2.
\]
Thus
\[
 x^2=(m^2-2n^2)(m^2+2n^2).
\]
The two positive odd factors are coprime, so each is a square:
\[
 m^2-2n^2=u^2,
 \qquad
 m^2+2n^2=v^2.
\]
Then
\[
 \frac{v-u}{2}\cdot\frac{v+u}{2}=n^2.
\]
The two factors are coprime, so write them as $b^2$ and $a^2$.  Hence
$u=a^2-b^2$, $n=ab$, and
\[
 m^2=u^2+2n^2=(a^2-b^2)^2+2a^2b^2=a^4+b^4.
\]
This is another positive solution of the same form with hypotenuse $m<z$, contradicting minimality.
\end{proof}

\begin{thm}[local obstruction principle]
If an integer equation $f(x,y,z)=0$ has an integer solution, then for every modulus $m\ge2$ the congruence
\[
 f(x,y,z)\equiv0\pmod m
\]
has a solution.  Thus a single modulus with no congruence solution proves that the integer equation has no solution.
\end{thm}

\begin{app}
\begin{enumerate}[label=(\alph*)]
\item The equation
\[
 x^3+y^3-7z^3=3
\]
has no integer solution.  Modulo $7$, every cube is $0$, $1$, or $-1$, so $x^3+y^3$ can only be $-2,-1,0,1,2$, never $3$.

\item The equation
\[
 y^2=x^3+7
\]
has no integer solution.  If $x$ is even, the right side is $3$ modulo $4$.  Hence $x$ is odd.  The case $x\equiv3\pmod4$ would make the right side $2$ modulo $4$, so $x\equiv1\pmod4$.  The equation also forces $x\ge1$.  Moreover,
\[
 y^2+1=x^3+8=(x+2)(x^2-2x+4).
\]
Since $x+2\equiv3\pmod4$, some prime $p\equiv3\pmod4$ divides $x+2$ to an odd power.  Thus $-1$ is a quadratic residue modulo $p$, contradicting the first supplementary law.
\end{enumerate}
\end{app}

\chapter{Roots of Unity and M\"obius Inversion}

\begin{de}
A complex number $z$ is an $n$th root of unity if $z^n=1$.  It is a primitive $n$th root of unity if it is not a root of unity of any smaller positive order.
\end{de}

\begin{thm}
If $\zeta$ is a primitive $n$th root of unity, then
\[
 \zeta,\zeta^2,\ldots,\zeta^n
\]
are precisely the $n$ distinct roots of $x^n-1$.
\end{thm}
\begin{proof}
Each power is an $n$th root.  If $\zeta^i=\zeta^j$, then
$\zeta^{i-j}=1$, so $n\mid i-j$; among $1,\ldots,n$ this forces $i=j$.  A degree-$n$ polynomial has at most $n$ roots.
\end{proof}

\begin{thm}
If $\zeta$ is a primitive $n$th root of unity, then $\zeta^k$ is primitive if and only if
\[
 \gcd(k,n)=1.
\]
\end{thm}
\begin{proof}
The order formula is
\[
 \ord(\zeta^k)=\frac n{\gcd(k,n)},
\]
proved exactly as in Theorem 4.5.
\end{proof}

\begin{de}
Let $F(n)$ be the sum of all $n$th roots of unity, and let $f(n)$ be the sum of all primitive $n$th roots of unity.
\end{de}

\begin{thm}
For every $n\in\N$,
\[
 F(n)=I(n)=
 \begin{cases}
 1,&n=1,\\
 0,&n>1,
 \end{cases}
 \qquad
 f(n)=\mobius(n).
\]
\end{thm}
\begin{proof}
For $n>1$, the roots of $x^n-1$ have sum $0$ by Vieta's formula; for $n=1$, the only root is $1$.  Every $n$th root has a unique exact order $d\mid n$, so
\[
 F(n)=\sum_{d\mid n}f(d).
\]
M\"obius inversion gives
$f(n)=\sum_{d\mid n}\mobius(d)F(n/d)=\mobius(n)$.
\end{proof}

\begin{de}
Let $\zeta_n$ be a primitive $n$th root of unity.  The $n$th cyclotomic polynomial is
\[
 \Phi_n(x)=\prod_{\substack{1\le k\le n\\\gcd(k,n)=1}}
 (x-\zeta_n^k).
\]
\end{de}

\begin{thm}
For every $n\in\N$,
\[
 \Phi_n(x)\in\Z[x].
\]
\end{thm}
\begin{proof}
The roots of $x^n-1$ split according to their exact orders, giving
\[
 x^n-1=\prod_{d\mid n}\Phi_d(x).
\]
Induct on $n$.  For every proper divisor $d<n$, the polynomial $\Phi_d$ is monic with integer coefficients.  Their product is therefore monic in $\Z[x]$ and divides $x^n-1$ in $\C[x]$.  Long division by a monic integer polynomial produces an integer quotient, which is $\Phi_n$.
\end{proof}

\begin{thm}[explicit formula]
For every $n\in\N$,
\[
 \Phi_n(x)=\prod_{d\mid n}(x^d-1)^{\mobius(n/d)}.
\]
\end{thm}
\begin{proof}
The factorization
$x^n-1=\prod_{d\mid n}\Phi_d(x)$
may be viewed in the multiplicative abelian group $\C(x)^\times$.  Applying M\"obius inversion to this divisor-product identity yields
\[
 \Phi_n(x)=\prod_{d\mid n}(x^d-1)^{\mobius(n/d)}.
\]
Although negative exponents appear in the formal product, Theorem 8.4 shows that the result simplifies to a polynomial in $\Z[x]$.
\end{proof}


\end{document}
